\documentclass[11pt,letterpaper]{article}

%% =====================================================================
%%  DART -- submission version
%%
%%  Placeholders that must be filled before submission are marked
%%  \TODO{...} and render as [[ ... ]] in the PDF so none can be missed.
%%  Set \submissionfalse to hide line numbers for a clean reading copy.
%% =====================================================================

%% ---- Packages ----
\usepackage[utf8]{inputenc}
\usepackage[T1]{fontenc}
\usepackage{mathpazo}
\usepackage[margin=2cm]{geometry}
\usepackage{graphicx}
\usepackage{float}
\usepackage{caption}
% Captions: single-spaced and one size down, as Nature sets them. The double-spacing the
% submission switch applies to body text is for reviewers to write between lines of PROSE; applying
% it to a 600-word figure legend pushes the float past a full page. singlelinecheck=false skips
% TeX's single-line measuring box, which overflows on captions this long.
\captionsetup{singlelinecheck=false, font={small, singlespacing}, skip=4pt}
\usepackage{amsmath,amssymb}
\usepackage{booktabs}
\usepackage[dvipsnames]{xcolor}
\usepackage[most]{tcolorbox}
\usepackage{hyperref}
\usepackage[numbers,sort&compress,super]{natbib}
\usepackage{enumitem}
\usepackage{lineno}
\usepackage{etoolbox}

%% ---- Submission switches ----
\newif\ifsubmission
\submissiontrue          % line numbers + double spacing for review
\ifsubmission
  \linenumbers
  \usepackage{setspace}
  \onehalfspacing
\fi

%% ---- Unfilled-content marker ----
\newcommand{\TODO}[1]{{\color{BrickRed}\textbf{[[#1]]}}}

%% Control drugs (above threshold at no recorded seed) that nonetheless cross the
%% +0.01 threshold under seed resampling. Single source of truth for the three
%% places this range is quoted; see the seed-stability table in the code release.
\newcommand{\CONTROLRANGE}{1 to 4}

%% The premise correlation, in two forms (analysis/natural/zhao_premise_disjoint.py;
%% results/zhao_gbm/premise_disjoint.csv). The first correlates the mean signature with the
%% malignant-compartment response, and the mean signature CONTAINS 43% of those cells, so part of
%% it is arithmetic. The second correlates two compartments that share no cells. Both are quoted;
%% the disjoint one is the load-bearing figure.
\newcommand{\PREMISEOVERLAP}{+0.878}
\newcommand{\PREMISEDISJOINT}{+0.835}

%% Class C, semantically matched functional oracle (analysis/class_c/class_c_functional_oracle.py;
%% results/upgrade/class_c_functional_oracle.json). Defined once so no number can drift between
%% the text, the figure caption and the Methods.
\newcommand{\CTXLINES}{966}
\newcommand{\NGDSCDRUGS}{286}
\newcommand{\NCQUERIES}{103}
\newcommand{\ENERGYFUNC}{+0.276}
\newcommand{\MEANCTRLFUNC}{+0.083}
\newcommand{\MEANRAWFUNC}{+0.241}
\newcommand{\MAGMATCHFUNC}{+0.232}
\newcommand{\MAGONLYFUNC}{+0.330}
\newcommand{\POTMATCHFUNC}{+0.399}
\newcommand{\ENERGYPARTIAL}{+0.097}
\newcommand{\ENERGYWINS}{62}
\newcommand{\ENERGYWINSP}{0.073}
\newcommand{\CENTERUNCEN}{85\% of drug pairs positive, median $+0.178$, versus 52\% and
  $+0.015$ once corrected}

%% ---- Section formatting ----
\usepackage{titlesec}
\titleformat{\section}{\large\bfseries}{}{0em}{}
\titleformat{\subsection}{\normalsize\bfseries}{}{0em}{}
\titlespacing*{\section}{0pt}{1.5ex plus 0.5ex}{0.8ex plus 0.2ex}
\titlespacing*{\subsection}{0pt}{1.2ex plus 0.3ex}{0.5ex plus 0.1ex}

\hypersetup{colorlinks=true,linkcolor=MidnightBlue,citecolor=MidnightBlue,urlcolor=MidnightBlue}
\setlength{\emergencystretch}{3em}

%% ---- Float placement ----
%% Without these, a figure that is tall relative to the text block gets promoted to a float page and
%% strands the rest of the sheet blank (Fig. 3 was doing exactly that). LaTeX's defaults reserve 20%
%% of every page for text (textfraction) and refuse to let a top float exceed 70% (topfraction), so a
%% 6.9 in figure plus a legend has nowhere legal to sit and is deferred. These values let a large
%% float share its page with running text instead.
\renewcommand{\topfraction}{0.92}
\renewcommand{\bottomfraction}{0.85}
\renewcommand{\textfraction}{0.06}
\renewcommand{\floatpagefraction}{0.65}
\setcounter{topnumber}{2}
\setcounter{bottomnumber}{1}
\setcounter{totalnumber}{3}

\begin{document}

%% =====================================================================
%%  TITLE PAGE
%% =====================================================================
\thispagestyle{empty}
\begin{flushleft}

{\LARGE\bfseries Objective-aligned evaluation inflates the distributional gain\\[2pt]
in single-cell drug retrieval\par}

\bigskip
\noindent\textbf{Running title:} Objective--utility mismatch in single-cell drug retrieval

\bigskip
\noindent\textbf{Authors}\\
\TODO{Author One}$^{1}$\TODO{ORCID}, \TODO{Author Two}$^{2}$\TODO{ORCID}, \TODO{additional authors}, \TODO{Senior Author}$^{1,}$\footnote{Correspondence: \TODO{email}}

\medskip
\noindent\textbf{Affiliations}\\
$^{1}$\TODO{Department, Institution, City, Country}\\
$^{2}$\TODO{Department, Institution, City, Country}

\medskip
\noindent\textbf{Corresponding author:} \TODO{Name}, \TODO{email}

\bigskip
\noindent\textbf{Keywords:} single-cell perturbation; drug retrieval; distributional distance;
objective--utility mismatch; benchmark design; virtual cell

\bigskip
\noindent\textbf{Word count:} abstract 150 words; main text (Introduction, Results, Discussion,
excluding figure captions) approximately 11{,}700 words; Methods 4{,}897 words. Display items: 6 main
figures, no main-text tables. Supplementary Information: 6 Extended Data figures, 3 Supplementary
Notes, 4 Supplementary Tables.

\end{flushleft}

\vfill
\noindent\rule{\textwidth}{0.4pt}

\clearpage

%% =====================================================================
%%  ABSTRACT
%% =====================================================================
\begin{abstract}
\noindent
Single-cell drug retrieval is moving from mean signatures to distributional scores over treated populations, and routinely reports large gains under a metric of its own kind. We audit that practice under increasingly independent criteria, and show that mean-signature retrieval is the distributional family's zero-variance limit, not a competitor. Holding rankings, scorer and queries fixed and changing only the judge, distributional retrieval is decisively ahead under objective-aligned metrics, reverses under mechanism recovery, and beats its incumbent against an external functional oracle, where a scalar comparing no distributions reproduces most of it. Even blind external oracles reverse the winner by statistical form. We formalise when cell-level structure can pay off as two gates, differential response and recoverability; in patient tumours nobody constructed both open, yet no advantage appears: the mean already ranks most of what a subpopulation does. Distributional information here is real, and much smaller than objective-aligned evaluation makes it appear.
\end{abstract}

\bigskip\noindent\rule{\textwidth}{0.4pt}\bigskip

%% =====================================================================
%%  INTRODUCTION
%% =====================================================================
\section*{Introduction}

Single-cell technologies have changed what it means to describe a drug's effect: it is not a single number but a distribution. Within a treated population, distinct subpopulations can respond in opposite directions, and the state that governs therapeutic outcome is often a minority one, a persister cell or a resistant clone. This premise drives the current wave of investment in virtual-cell models\cite{AIVirtualCell}, single-cell foundation models\cite{ref7,ref8,ref9}, and perturbation-response predictors\cite{ref4,ref5,ref10}: if responses can be characterized at single-cell resolution, one should be able to rank candidate drugs better than methods that average the population into a mean signature, as the Connectivity Map does\cite{ref1,ref2}.

That premise is rapidly becoming method. Recent work has brought population-to-population distributional distances, optimal transport, maximum-mean-discrepancy, energy distance and the Wasserstein metric, into single-cell perturbation modeling, and these methods commonly report substantial gains over mean-based approaches. Yet they are almost always scored by distributional metrics of their own kind. When a distributional model is graded by a distributional yardstick, how much of the reported gain is a genuine improvement in decisions, and how much is the method and the yardstick sharing a mathematical target, has been asked too rarely\cite{ref3,ref13,ref15}, though the general hazard of steering by a proxy is not new\cite{MetricGaming}. We name this an objective--utility mismatch, and define it narrowly: \emph{a retrieval system is evaluated by a criterion that measures the same statistical property its ranking maximised, so fidelity to the objective is read as evidence of downstream utility}. The term is scoped to retrieval evaluation and we do not extend it beyond that here. We are equally deliberate about how strong a word it licenses. Training on an objective and testing with a metric of the same family is not, by itself, a defect: a classifier fit by cross-entropy and scored by held-out log loss is honestly evaluated, because held-out data still tests generalization and the model can still lose. Retrieval is different, and sharper. There is no fitting step and no held-out set; the candidate is chosen by \emph{maximizing a score}, and the evaluation then measures \emph{that same score's} ground truth on \emph{those same candidates}, so the ranking function and the evaluation function are one function evaluated once, with no estimation error to expose and no way to lose. We reserve the word circular for exactly this condition, and state it formally in Results.

Deciding whether ``distributional methods are better'' survives past the field's own metrics requires an oracle that is both independent of the score and \emph{semantically matched to the retrieval task}; getting the second requirement wrong is its own trap, since a similarity retriever graded on absolute potency is being asked a question it was never posed. The distinction we hold throughout is between objective fidelity and independent utility (Fig.~\ref{fig:1}). Objective-aligned metrics (Class A) legitimately test whether a method computes its own objective; task-proximal metrics such as mechanism-of-action recovery (Class B) are more independent but remain derived from the same upstream data; oracle-independent utility metrics (Class C), such as dose-response viability or longitudinal survival, are what a therapeutic claim ultimately requires, and are largely absent from the benchmarks these methods are evaluated on. An audit across six published method families, spanning signature retrieval\cite{ref1,ref2}, expression prediction\cite{ref4,ref5,ref6} and graph-based inverse design\cite{ref10,ref11,ref12}, finds objective-aligned evaluation to be the default in each. The central question is therefore not whether distribution-aware retrieval beats mean retrieval, but where objective fidelity ends and independent utility begins.

Making that boundary measurable calls for an instrument rather than a proposal: one family of population-to-population retrieval scores whose candidate rankings can be re-graded under criteria of increasing independence without being recomputed. We implement that family as DART (Distributional Auditing of Retrieval Transfer) and release it with the code, and it is used here as a probe rather than offered as a method. Mean-signature retrieval, including Connectivity-Map cosine, is not a competitor but the zero-variance collapsed member of this same family: under the implemented kernel the two are numerically identical, with Hit@1 agreeing to three decimals. That unification is what makes the probe a controlled one, since only the judge then changes while the rankings, the scorer and the queries stay fixed, and we watch where the conclusion moves. Under objective-aligned metrics distributional retrieval looks decisively ahead; re-scored for mechanism recovery on the identical rankings the lead does not merely shrink but reverses; against an external functional oracle it does beat its incumbent, yet a scalar that compares no distributions reproduces almost the whole margin. The mismatch even reaches past metrics a method can see: two oracles built from the same proteins in the same cells, differing only in whether each is computed as a mean or as a distribution, hand victory to opposite methods.

What bounds the distributional gain is mechanistic, and it is the affirmative payload of this work. A predictor that adds a \emph{single} delta vector to every cell has, by algebra rather than by experiment, two subpopulations with identical responses, so a distributional score has nothing to exploit; usefulness therefore requires that subpopulations both respond \emph{differently} and can be \emph{told apart} at all. We formalize this as a two-gate criterion, differential response and recoverability, and test it where it matters: in patient-derived glioblastoma tumours that nobody constructed, rather than in the cell-line mixtures we had assembled, which turned out to distort the biology in both directions at once and to flatter the premise both times. In real tissue both gates open, and the advantage still fails to appear, because the mean already ranks most of what the compartments do. The two gates are necessary, not sufficient, and what they omit is a third condition we name but do not measure as we measure them: the differential response must change \emph{which candidate is preferred}, and in this tumour it largely does not.

Our angle differs from two adjacent lines of work. One provides distributional \emph{prediction} methods, using optimal-transport or Schr\"odinger-bridge formulations to model perturbation responses\cite{CellOT,Departures}; these ask how to predict outcomes better, not how to rank candidate drugs, and do not test whether the evaluation is circular. A second compares distributional evaluation metrics for generative models\cite{OptimalMetrics,StandardizedEval}, asking which metric better measures whether a generated distribution matches a real one; we ask the orthogonal question of how much of a reported gain is self-fulfilling when the retrieval score and the evaluation metric are the same object. To our knowledge neither this circularity nor the zero-variance identity has been stated before.

Three contributions follow, in the order the Results establish them. First, mean and distributional retrieval form one continuous spectrum of variance sensitivity rather than competing method classes, with mean-signature retrieval as its zero-variance endpoint; that is algebra, and it is what licenses holding the scorer fixed while the judge changes. Second, how large a distributional gain looks depends on how independent that judge is: decisive under an objective-aligned metric, reversed under mechanism recovery, and, for one pair of external oracles, settled by the oracle's statistical form alone. Third, independent validation leaves a distributional contribution that is smaller than reported rather than absent, since against a semantically matched functional oracle the distributional score does beat its incumbent while a scalar comparing no distributions reproduces most of that margin. What this offers the field is a test for deciding whether a distributional gain is real or circular, a diagnosis of a failure mode that likely inflates results well beyond drug retrieval\cite{ref16,ref17}, and a roadmap of measurable conditions under which cell-level structure can pay off. The roadmap points somewhere specific, and it is the one encouraging thing we found: in a real tumour the subpopulation information a distributional score would need is \emph{present}, and it is unsupervised clustering, not the data, that fails to recover it. That is an engineering problem, and engineering problems get solved. As single-cell foundation models proliferate, the question will only become more pressing.

%% =====================================================================
%%  RESULTS
%% =====================================================================
\section*{Results}

Separating objective fidelity from independent utility requires an apparatus that varies the
evaluator while holding everything else fixed. The apparatus is a single family of
population-to-population retrieval scores, implemented as DART, whose candidate rankings can be
re-graded by criteria of increasing independence without recomputing the rankings themselves
(Fig.~\ref{fig:1}). The same
candidates, ranked once, are scored by the retrieval score's own distributional objective (Class A), by a
task-proximal biological criterion such as mechanism recovery (Class B), and by an
oracle-independent functional readout (Class C), so that from one analysis to the next the
rankings, the scorer and the queries are held fixed and only the evaluator changes. Each apparent
advantage thereby becomes a testable claim: a distributional gain that reflects better decisions
should survive a more independent judge, whereas one that reflects only the scoring of a method
with a metric of its own kind should not.

% FIGURE 1
\begin{figure}[!htbp]
\centering
\includegraphics[width=\textwidth]{figures/fig1.pdf}
\caption{\textbf{Distributional differences are visible, but their value depends on the evaluator.} \textbf{a}, The premise, drawn on one response coordinate: two drugs produce an identical mean shift, but under drug B the resistant minority does not move and the responding bulk overshoots to compensate, so the population mean is preserved while the two drugs do opposite things to the cells that matter. Schematic. \textbf{b}, The inverse retrieval task and the two ways to score it: a distributional score keeps the whole candidate population, whereas mean-signature cosine collapses it to a single delta vector, and both return one ranked candidate list. \textbf{c}, Two candidate populations placed on the same sample mean, with the marginal densities above the scatter showing that one is bimodal and the other unimodal: means tie, distributions separate. \textbf{d}, When an evaluator is independent. A retrieval score graded by a metric of its own kind is one function evaluated once (coupled, Class A), so the ranking cannot lose; the same ranking sent to an outside judge can be confirmed or overturned (Class B/C). \textbf{e}, The three evidence classes ordered by independence from the retrieval objective, with the headline metric of published method families (filled for Class A only) against what this study reports (all three, its Class-C readout imported). Its Class-C readout is an imported, semi-real functional oracle built from dose-response profiles in cell lines the transcriptional data never saw; a Class-C claim measured on the same cells would require wet-lab validation.}
\label{fig:1}
\end{figure}

\subsection*{Mean and distributional retrieval are one continuum of variance sensitivity}

We define the inverse retrieval task at the level of populations. For a target
population $Q$ and a library of candidate perturbation populations $\{P_d\}$, a
retrieval score $s(P_d, Q)$ ranks candidates by similarity of effect. Mean-signature
retrieval replaces each population with its mean differential-expression vector and
scores the cosine (or correlation) between mean signatures, the CMap operation.
Distributional retrieval instead scores the full populations with a distributional distance (energy,
MMD, sliced-Wasserstein, or a subpopulation-coverage aggregate) and ranks by the
negative distance.

These are not two unrelated method classes but two ends of one spectrum. Writing a
population as its mean plus a zero-mean residual, the distributional distance
reduces exactly to the mean-to-mean distance when the residual variance is taken to
zero. Mean-signature retrieval is therefore the zero-variance collapsed member
of the distributional family. Under the implemented cosine kernel this is not an approximation
but an identity: mean-cosine retrieval and CMap cosine retrieval achieve identical Hit@1
(0.3885, to six figures 0.388492), because they compute the same operation, the cosine
of mean-delta signatures. Distributional energy retrieval reaches 0.837 on the same
queries. Both figures are unweighted macro-means over the seven (task $\times$ setting) cells of
three retrieval tasks: controlled SciPlex3, cross-line SciPlex3 and Frangieh.

The collapse is one end of a continuous spectrum. The subpopulation-coverage aggregate carries
a temperature $\beta$ that interpolates between aggregation regimes, and on the controlled
benchmark the interpolated distance rises monotonically from 0.7125 at $\beta \to 0$,
exactly the mean-aggregated value, to 1.300 at $\beta \to \infty$, exactly the worst-case
value (Fig.~\ref{fig:2}e). The global energy distance sits inside the same family, equalling the
single-partition ($K = 1$) coverage aggregate exactly (Panobinostat, energy $=$ coverage$_{K=1}$
$=$ 0.1099, difference 0). Energy, coverage and mean retrieval are therefore one axis of variance
sensitivity rather than competing method classes, and mean-signature retrieval is its
zero-variance endpoint (Fig.~\ref{fig:2}c,f).

This unification holds independently of any downstream utility metric: it is an
algebraic fact about the scoring rules, not a claim about which method is
therapeutically better. Two scope limits are explicit. First, the exact identity is
specific to the cosine connectivity kernel; the canonical L1000 WTCS score is a
rank-based (GSEA-style) enrichment of the same mean-collapsed signature, monotone in the
same quantity but not algebraically identical to the cosine kernel implemented here (cmap\_wtcs Hit@1
0.4635, which differs). Second, the identity concerns the scoring operation, not the
choice of input population; the value of retaining the full distribution is exactly what
the remaining results interrogate.

% FIGURE 2
\begin{figure}[!htbp]
\centering
\includegraphics[width=\textwidth]{figures/fig2.pdf}
\caption{\textbf{Mean and distributional retrieval form a continuous spectrum of variance sensitivity.} \textbf{a}, Generative model $X=\mu+\lambda\epsilon$: the mean signature is the $\lambda\to0$ limit. \textbf{b}, Scaling only residual variance leaves the mean fixed. \textbf{c}, Energy distance converges to the mean-distance floor as residual spread vanishes. \textbf{d}, Under the implemented cosine kernel, mean cosine and CMap cosine are numerically identical (0.388); WTCS is a rank-based sibling. \textbf{e}, The $\beta$-interpolation spectrum from mean-aggregation to worst-case, with energy equal to single-cluster coverage. \textbf{f}, Metric-family hierarchy: every limit collapses to the mean.}
\label{fig:2}
\end{figure}

\subsection*{Objective-aligned evaluation inflates the gain}

The collapse identity above has an empirical face. Across all retrieval metrics, mean cosine and
Connectivity-Map cosine are perfectly correlated ($\rho = 1.00$), while the energy and coverage
scores form a separate, near-orthogonal block (Extended Data Fig.~1, which also gives dataset scale
and the distributional metric definitions). The correlation is computed over the 54,180
query--candidate scores that the 1,260 retrieval queries generate against their candidate
libraries, at 43 candidates per query. The unit is the scored pair, not the query: an earlier draft
called these ``54,180 real queries'', which overstates the independent sample size by a factor
of 43.

When candidate responses are \textit{observed}, so that full candidate populations are
available rather than predicted, distributional retrieval looks decisively stronger than mean
retrieval. Energy-distance retrieval reaches Hit@1 0.837 against 0.388 for
mean-cosine and the numerically identical CMap cosine; PCA-distance retrieval reaches 0.778 against
PCA-mean's 0.518 (Fig.~\ref{fig:3}). These are unweighted macro-means over the seven (task
$\times$ setting) cells of the three retrieval tasks. Weighting cells by their query counts gives
0.887 against 0.421; we report the unweighted figure throughout, because it does not let the
largest cell dominate.

Extending from observed populations to a partial-observed retrieval setting, we score candidates
by an energy-based welfare regret, how much decision regret a ranking incurs against a latent
welfare objective. Over all 765 partial-observed queries, the worst-case
subpopulation-coverage scorer reduces median regret by \textbf{$+$0.118} against mean retrieval
(mean $+$0.258), improving 72\% of queries (paired Wilcoxon
$p = 1.6\times10^{-66}$). We report the full query set here rather than the
gate-recommended subset of it: the two are indistinguishable ($+$0.119 on the 621 recommended
queries), and selecting a headline on a gate-selected subset would be exactly the analytic degree
of freedom this paper audits.

By the standard of these metrics, distributional retrieval is far ahead. But the metric
here is objective-aligned: the energy welfare regret and the energy retrieval score measure
the same distributional object. This result establishes objective fidelity: distributional retrieval computes a
discriminative objective, and observed heterogeneous populations contain real distributional
signal. It is a \textit{distributional retrieval signal under an objective-aligned proxy}, not
evidence of independent therapeutic utility. Whether that signal survives a more independent
judge is the question the next sections answer.

% FIGURE 3
\begin{figure}[!htbp]
\centering
\includegraphics[width=\textwidth]{figures/fig3.pdf}
\caption{\textbf{Under objective-aligned metrics, distributional retrieval looks decisively stronger.} \textbf{a}, Hit@1 ladder: energy 0.837 down to mean/CMap cosine 0.388, unweighted macro-mean over the seven (task $\times$ setting) cells of three retrieval tasks. \textbf{b}, Per-task Hit@1. The advantage holds on the controlled and cross-line SciPlex3 tasks but \emph{not} on Frangieh, the only natural real dataset in the macro-mean, where mean/CMap cosine (0.600) beats energy (0.578). The macro-mean hides this counterexample; we state it here and in the text. \textbf{c}, Per-query regret reduction over \emph{all} 765 partial-observed queries: median $+0.118$, 72\% improved (the $+0.119$ quoted elsewhere is the 621-query gate-recommended subset, and the two are indistinguishable, which is panel \textbf{d}'s point). \textbf{d}, Recommended ($+0.119$) versus non-recommended ($+0.122$) subsets: the gate does not concentrate the gain. \textbf{e}, Alpha-crossover. The energy advantage does not close monotonically as the two subpopulations merge, and it does so in only one of the three cell lines (K562, from $1.00$ to $0.05$ in Hit@1 as $\alpha$ rises from $0.5$ to $0.9$); in A549 and MCF7 the gap is as large or larger at high $\alpha$. We report the per-line curves rather than a pooled trend for that reason. \textbf{f}, All five distributional metrics show positive Class-A regret reduction.}
\label{fig:3}
\end{figure}

\subsection*{Task-proximal metrics collapse the apparent gains}

This is the central result (Fig.~\ref{fig:4}). Both sides of the comparison use one
scorer on one query set, so the only thing that changes is the class of metric doing the
judging. Selecting the best scorer separately for each metric class, on different query subsets, is
an analytic degree of freedom this task readily permits, and it would report $+0.119$ against
$-0.013$: it \emph{understates} the collapse. We fix both and report the larger number.

On the 600 partial-observed queries where both metric classes are defined (MoA recovery is
undefined for the remaining 165; Methods), the worst-case coverage scorer achieves a Class-A
regret reduction of $+$0.129 (median; mean $+$0.275). Re-scoring the \emph{identical}
rankings under mechanism-of-action recovery nDCG, the same two statistics are 0.000
(median) and $-$0.037 (mean; Fig.~\ref{fig:4}a). We give both, because they say
different things and both are load-bearing: on the median query the advantage does not shrink
but vanishes \emph{exactly}, and averaged over queries it turns \emph{negative}. The reversal is
not an artefact of a heavy tail: 41.3\% of queries \emph{lose} mechanism recovery under the
distributional score against 34.2\% that gain it (Wilcoxon $p = 2.4\times10^{-4}$). The distributional score
is not merely no better; it is worse. All five distributional variants behave the same way (Class-A gain $+$0.056
to $+$0.129, Class-B loss $-$0.011 to $-$0.037), so the conclusion does not depend on which
variant is chosen.

Projecting the real datasets onto the retrieval task, the distributional advantage in
minority-state coverage across all 239 real-data tasks is statistically real and
practically negligible: mean $+$0.0018, Wilcoxon $p = 1.3\times10^{-7}$ on the
137 tasks with a nonzero difference, on a metric whose observed range is $[0.60, 1.00]$.
In 102 of the 239 tasks the difference is \emph{exactly} zero, because both method
families select the same drug. And the effect is confined to the mixtures we
constructed: the mean gain is $+$0.0043 on the 90 constructed cross-line tasks, against
$+$0.0004 within-line ($n = 90$), $+$0.0004 on Frangieh ($n = 23$), $+$0.0003 on CD34$+$
($n = 12$) and $+$0.00003 under predicted candidates ($n = 24$) (Fig.~\ref{fig:4}g).

An independent divergence audit reaches the same verdict and numerically the
same effect: a coverage gain that is real, monotone in divergence, and around $+$0.002,
which is to say too small to act on. Two analyses that share no code path converge on it.

\textbf{A dominance count on this data counts noise excursions.} The conventional summary would
count tasks whose gain clears a threshold, here $+0.01$. That count is not interpretable
(Methods): the threshold sits inside the noise band of the difference distribution and is
one-sided, and the seed is not a replicate, since it re-draws which drugs are held out and
re-clusters the minority subpopulation. Holding the drug fixed and resampling the seed ten times,
only \textbf{2 of the 10} threshold-crossing cross-line tasks stay above it in at least 8 of 10
seeds, while control drugs that clear it at \emph{no} recorded seed clear it in \CONTROLRANGE{} of
10 resamples. We report the distribution instead. The two survivors are a hypothesis rather than a
result: A549$\rightarrow$MCF7 Abexinostat ($+0.090 \pm 0.013$) and Belinostat
($+0.081 \pm 0.027$), both HDAC inhibitors, both above threshold in 10 of 10 seeds. With $n = 2$,
one cell-line pair and one drug class, this is the only place in this study where a distributional
advantage on real data survives being checked, and it is suggestive and no more.

Crucially, a diagnostic gate designed to \textit{concentrate} the advantage onto the
queries where distributional information should help does not do so. The median
energy-proxy regret reduction is $+$0.119 on the gate-recommended subset ($n = 621$)
against $+$0.122 on the non-recommended subset ($n = 133$; the remaining 11 of the 765
queries are labelled mean-sufficient). The gate does not separate the cases where the
distributional score helps from those where it does not, even on the proxy metric it was tuned
against, and if anything it selects slightly against them.

The simplest demonstration is at the dataset level. The divergence criterion holds that
distributional retrieval can help only when subpopulation responses diverge,
operationalized as a response cosine below $\approx$0.9. Measuring where each dataset
actually sits on that axis, cross-line SciPlex3 lies at cosine 0.032, Frangieh
at 0.709, and CD34$+$ at 0.186. All three clear the gate. Yet the distributional score wins
only on cross-line: on Frangieh, the only natural real dataset, the mean-signature
incumbent wins (Hit@1 0.600 versus 0.578), and CD34$+$ is our \emph{negative control},
where the mean wins decisively. A criterion that admits the positive anchor and
the negative control alike is not a criterion, and the dataset-level divergence rule
therefore fails on its own terms.

A direct audit makes the failure concrete (Fig.~\ref{fig:4}e,f). The gate's
structure-reliability axis is anti-correlated with true response divergence (Spearman
$\rho = -0.21$, $p = 3.8\times10^{-9}$), moving opposite to the quantity it is meant to track;
its preference-conflict axis is only weakly aligned ($\rho = +0.18$); and the recommendation
label does not separate high- from low-divergence queries (Mann--Whitney $p = 0.090$).
Bypassing the gate and stratifying directly by true divergence, the highest-divergence stratum
shows a statistically significant but negligible minority-coverage gain (median gap rising from
$+$0.0009 in the lowest quartile to $+$0.0018 in the highest, all Benjamini--Hochberg
$q \approx 0$, against a metric operating range of $[0.89, 0.99]$: about 0.13 standard
deviations). Mechanism recovery is null in the highest stratum ($q = 0.57$, $n = 143$) and
\emph{significantly negative} in the lowest (median gap $-$0.030, $q = 5.8\times10^{-5}$,
$n = 166$): where responses are least divergent, the distributional score actively degrades annotation
recovery.
The deficit is in the data and the task, not only in the gate.

The obvious objection is that a null is merely an underpowered study, and the two metric
classes are indeed unequally powered (Extended Data Fig.~2): the coverage metric is fully
powered (achieved power 1.00), while MoA recovery is not (0.06; 80\% power would need roughly
20,800 queries against the 143 available). We therefore do \emph{not} claim to have proven the
absence of a small MoA effect. We claim something weaker and sufficient: where the theory
predicts the largest gain there is no detectable gain, and where responses are least divergent
the effect that \emph{is} detectable runs against the distributional score.

The conclusion is deliberately strong and deliberately bounded: apparent
distributional gains are metric-dependent and cannot be interpreted as therapeutic
utility without independent validation. One scorer, one query set and one set of
rankings is ``far ahead'' or ``worse than the incumbent'' depending only on which metric
class judges it. That dependence, not a verdict on any single method, is the finding.

% FIGURE 4
\begin{figure}[!htbp]
\centering
\includegraphics[width=\textwidth]{figures/fig4.pdf}
\caption{\textbf{Independent evaluation collapses the apparent advantage, and what survives it is mostly magnitude.} The metric-class ladder, read row by row.
\textbf{Row 1, Class A $\rightarrow$ Class B.} \textbf{a}, One scorer, one query set, only the class of metric changes: a Class-A regret reduction of $+0.129$ becomes a Class-B MoA-nDCG gain of $-0.037$. \textbf{b}, MoA-nDCG gain by cell line, uniformly at or below zero. \textbf{c}, Minority-coverage gain by divergence quartile: significant (all Benjamini--Hochberg $q \approx 0$) and negligible, every bar $\leq 0.006$ on a metric whose values run 0.89 to 0.99. Bars are quartile means with s.e.m.; dashes mark the medians the text quotes. \textbf{d}, The gate does not concentrate the advantage ($+0.119$ on 621 recommended queries versus $+0.122$ on 133 non-recommended).
\textbf{Row 2, why the gate cannot rescue it.} \textbf{e}, The gate's structure-reliability axis is \emph{anti}-correlated with true response divergence ($\rho=-0.21$), moving opposite to the quantity it is meant to track. \textbf{f}, Recommended and non-recommended queries have indistinguishable divergence (Mann--Whitney $p=0.09$). \textbf{g}, The real-data coverage advantage, one dot per task across all 239 tasks: real but negligible (overall mean $+0.0018$) and confined to the constructed cell-line mixtures ($+0.0043$ against $+0.0003$ in the natural data and $+0.00003$ under predicted candidates). We plot the difference rather than a dominance count; only 2 of the 10 threshold-crossing cross-line tasks survive seed resampling (circled).
\textbf{Row 3, Class C, an independent functional oracle.} \NCQUERIES\ leave-one-drug-out queries, SciPlex3 at 10\,$\mu$M. The oracle is \emph{drug--drug functional similarity} measured out of context, not absolute potency (Methods). \textbf{h}, Energy reaches $\rho=\ENERGYFUNC$ against $\MEANCTRLFUNC$ for the correctly specified mean incumbent, the first oracle-invisible criterion here on which the distributional score wins. It is not a \emph{distributional} win: a query-dependent scalar matching only response \emph{size} reaches $\MAGMATCHFUNC$ (hatched), and partialling that channel out leaves energy with $\ENERGYPARTIAL$ (dotted). The longest bar, potency matching at $\POTMATCHFUNC$, is a diagnostic rather than a recommender, since it needs the query drug's own measured potency. Open circles, per-cell-line medians. \textbf{i}, Per query, energy against that scalar: energy wins $\ENERGYWINS$ of \NCQUERIES\ queries, and the nominal $p=\ENERGYWINSP$ is anticonservative because queries within a line share a candidate pool. The mismatched absolute-potency comparison is reported only as a confounder audit (Extended Data Fig.~5).}
\label{fig:4}
\end{figure}

\subsection*{A functional oracle shows the gain is mostly magnitude}

Class B is task-proximal, not independent, so the ladder is incomplete without an external
functional readout. Choosing that readout is where a Class-C analysis is most easily got wrong,
and the wrong choice is the tempting one.

\textbf{The oracle must match the task.} A retrieval score is asked which candidate's response
distribution most resembles the query's. It is \emph{not} asked which candidate is most
cytotoxic. Grading a similarity retriever on absolute potency therefore tests a question it was
never posed: handed a weak query, a correctly working retriever \emph{should} return other weak
drugs. We report that mismatched comparison only as a confounder audit (below), because the
conclusion it invites, that retrieval is anti-aligned with therapeutic value, does not follow
from it.

The oracle that does match the task is drug--drug functional similarity: do the query
and the candidate \emph{behave} alike in an independent assay? GDSC2\cite{GDSC,GDSCLandscape} supplies this directly.
Each compound has a dose-response AUC profile across roughly a thousand cell lines, and two
drugs with similar cross-cell-line sensitivity patterns are functionally similar however potent
either happens to be. We define
$\mathrm{oracle}(q,c) = \mathrm{Spearman}\big(\mathrm{AUC}_{\bullet}(q),\,
\mathrm{AUC}_{\bullet}(c)\big)$ over the \CTXLINES\ GDSC2 cell lines that remain after
\emph{excluding} the three SciPlex3 lines, so the oracle is never evaluated in the context the
transcriptional query came from. Because a correlation is centered per drug, a compound's
overall potency level cannot drive it; only its selectivity \emph{pattern} can. Cell lines
differ enormously in general drug sensitivity, and left uncorrected that shared main effect makes
almost every drug pair look alike (\CENTERUNCEN); we therefore subtract each cell line's mean AUC,
estimated across all \NGDSCDRUGS\ GDSC2 compounds rather than only the matched ones, before
correlating.

\textbf{That correction decides the result, and we report the dependence rather than bury it.}
Without it, energy falls to $+0.097$ and drops \emph{below} the mean-signature incumbent's
$+0.138$, reversing the one affirmative finding of this section; the magnitude scalar
($+0.217$) beats both. Both oracles are given in full in Supplementary Table 4. The correction is
defensible on its face, it was specified on the reasoning above before the retrieval results were
inspected, and leaving a shared cell-line main effect in an oracle would be an error of its own.
But it is an analytic choice, and a paper arguing that the choice of criterion decides the winner
cannot exempt its own criterion from that finding. The honest statement is that the Class-C
positive below holds under the oracle we argue is correct and not under the oracle we argue is
wrong.

On this oracle, distributional retrieval beats its incumbent (Fig.~\ref{fig:4}h). Across
\NCQUERIES\ leave-one-drug-out queries, energy retrieval attains $\rho = \ENERGYFUNC$, against
$\MEANCTRLFUNC$ for the correctly specified mean-signature incumbent: a threefold margin, in the
same direction in all three cell lines. This is the first point in this study at which the
distributional score beats its incumbent on a criterion no scorer can see, and we report it as
such.

\textbf{It is nevertheless not a distributional result.} Rank the candidates instead by how
closely their response \emph{magnitude} matches the query's, a query-dependent scalar that
compares no distributions at all, and it reaches $\rho = \MAGMATCHFUNC$. Energy's edge over that
scalar is $\ENERGYWINS$ of \NCQUERIES\ queries, which even under an anticonservative paired test
is not significant ($p = \ENERGYWINSP$). Partial out the magnitude channel and energy retains
$\ENERGYPARTIAL$ (per line $+0.131$, $+0.053$, $+0.162$: positive in all three, but small).

Even a
\emph{query-independent} ordering carries comparable alignment: ranking candidates by their own
response magnitude, an ordering identical for every query that performs no retrieval whatsoever,
gives $\rho = -0.330$, so the same fixed ordering reversed gives $\MAGONLYFUNC$. We report its
magnitude rather than calling it a better recommender, and deliberately so: its \emph{sign} would
have to be known in advance, and a paired comparison in fact favours energy on 65 of \NCQUERIES\
queries even though its marginal median is the larger of the two. The point is only that a
constant ordering is not far behind a distributional score on this oracle. Finally, a single
scalar matching the candidate's potency to the query's reaches $\POTMATCHFUNC$, above the
distributional score, though it is a diagnostic rather than a usable recommender because it
requires the query drug's own measured potency; what it establishes is that the functional oracle
retains a substantial potency-level axis, which is the channel the mismatched comparison below
confounds entirely.

The picture is therefore layered rather than negative. Distributional retrieval carries genuine
functional signal that the mean signature does not; almost all of that signal is a
response-magnitude channel that a scalar reproduces; and the residual attributable to the
\emph{distribution} is roughly $\ENERGYPARTIAL$, about a third of energy's raw $\ENERGYFUNC$ and
positive but small. We deliberately do not express this as a ratio against the Class-A gain: a
Spearman correlation against an external oracle and a Hit@1 or regret improvement against an
objective-aligned proxy are different quantities on different scales, and dividing one by the
other would manufacture a number rather than measure one. The comparison that is legitimate is
within the oracle, and there two thirds of the distributional score's advantage is magnitude.
Objective fidelity does not make the distributional gain imaginary. It makes it look far larger
than it is.

One operational lesson is specific to baselines, and the other is deferred to the Discussion,
where we state the controls this study recommends. A mean-signature baseline scored \emph{without}
control subtraction reaches $\MEANRAWFUNC$, close to energy and far above the correctly specified
incumbent's $\MEANCTRLFUNC$, because it measures baseline cell identity rather than drug response.
A study using it as its incumbent compares against a broken baseline and will overstate its own
margin.

\textbf{Confounder audit: why the mismatched oracle inverts.} Scored against \emph{absolute}
potency, energy retrieval is anti-correlated ($\rho = -0.520$, same sign in all three lines) while
the query-independent magnitude scalar reaches $+0.692$. This is arithmetic, not pathology.
Response magnitude predicts potency on its own, and the energy \emph{distance} between query and
candidate largely tracks the candidate's own magnitude ($\rho = +0.791$), because big-response
populations sit far from everything, so ranking nearest-first returns the weakest candidates. The
same magnitude channel drives both oracles; this is not evidence that retrieval opposes therapeutic
value, and we do not claim that it is.

\textbf{Scope.} This oracle is semi-real: viability is measured in a separate bulk assay, joined by
compound identity rather than shared cells, under a different dose regime, so it constrains what a
transcriptional score recovers about drug function without being a clinical endpoint. The
\NCQUERIES\ queries within a cell line share a candidate pool, so paired tests across them are
anticonservative; the honest unit is the three cell lines, and we give per-line values throughout.
The strongest limit is the one stated above, that the affirmative result is not robust to the
cell-line centering: under the uncorrected oracle energy reaches only $+0.097$, below both the
incumbent ($+0.138$) and the magnitude scalar ($+0.217$). What survives both constructions is the
negative, that under neither oracle does the distributional score beat the scalar that compares no
distributions.

\subsection*{Oracle shape not biology decides which method wins}

The Class-C result above is not the only external readout available, and the second one disagrees
with it. Frangieh Perturb-CITE-seq measures 20 surface proteins in the \emph{same} 218,331 cells
whose RNA the scorers read, giving an oracle that is independent of the transcriptome, invisible to
every scorer, and free of the cross-assay join the GDSC oracle needs. Scored against it, the
\emph{mean} beats energy ($+0.242$ against $+0.146$), in all three immune conditions. Taken at face
value, one external oracle says distributional retrieval wins and another says it loses.

They are not in conflict. Look at what the protein oracle is: $\cos$ of the two knockouts'
\emph{mean} protein deltas. It collapses the protein readout to a mean, and the mean-cosine RNA
scorer computes that same statistic in the other modality. A mean-shaped oracle is structurally
matched to a mean-shaped scorer, and energy is being graded on a criterion that discarded the very
thing it compares. That is the mirror image of the Class-A alignment this paper is about, arising
now between two \emph{external} criteria, and it is testable directly.

\textbf{We therefore rebuilt the oracle from the same cells and the same proteins, twice, changing
only its statistical form} (Fig.~\ref{fig:5}c; Methods). One version is the mean cosine above; the
other is the energy distance between the two knockouts' protein-response \emph{cell clouds}, both
referred to that condition's control cells so it measures response and not baseline identity. The
two RNA rankings under test are byte-identical across the two columns. Nothing about the scorers,
the cells, the proteins or the queries changes. Only the shape of the judge changes.

\textbf{The winner swaps.} Under the mean-shaped oracle the incumbent leads ($+0.242$ against
$+0.146$); under the distribution-shaped oracle the distributional score leads, and by more
($+0.529$ against $+0.334$). The reversal holds independently in each of the three immune
conditions (Control, IFN$\gamma$ and Co-culture; $n = 213$, 227 and 219 queries). No scorer can see
either oracle, so this is not circularity in the Class-A sense of a method grading its own
objective. It is something we did not anticipate and consider the sharpest result in this paper:
even when the criterion is external, independent and blind, its own statistical form
selects the winner. A field that validates distributional methods against distributional readouts,
and mean-based methods against mean readouts, will find each vindicated, and neither finding will be
about biology.

\textbf{The magnitude control survives this, and refines it.} An energy distance tracks a
candidate's own response magnitude ($\rho = +0.791$), and both distributional objects here are
energy distances, so a magnitude-to-magnitude channel could in principle produce the swap with no
distribution ever being compared. It does not, quite. The magnitude scalar does climb steeply when
the oracle turns distributional (from $+0.125$ to $+0.352$), which is the confound behaving exactly
as predicted and is itself worth reporting. But energy still leads that scalar by $+0.177$ under the
distributional oracle, against only $+0.022$ under the mean-shaped one. So there is a genuine
distributional signal in these cells, and \emph{it is visible only to a criterion that is itself
distributional}. Whether that makes it a real advantage or an unfalsifiable one is the question this
paper leaves open, and we prefer to leave it open than to answer it with a criterion of our own
choosing.

\begin{figure}[!htbp]
\centering
\includegraphics[width=\textwidth]{figures/fig5.pdf}
\caption{\textbf{The statistical form of the criterion, not the biology, decides which method wins; and in real tissue both gates open without an advantage following.}
\textbf{Row 1.} \textbf{a}, HIR-Bench's analytic flip boundary $\alpha^*=B/(A+B)$ separates the regime in which the mean is a sufficient statistic from the one in which subpopulation structure can change the decision. \textbf{b}, Circularity measured rather than argued: features a method can observe at query time reach AUC 0.788 against a 0.643 majority-class rate, whereas features derived from the benchmark's own latent utility matrix, of which the label is also a function, are at chance (0.400). Pseudo-replication inflates the circular set by $+0.240$ and the honest one by only $+0.048$. Effective $n$ is 28 parameter cells, not 13,440 (Methods). \textbf{c}, The same effect between two external criteria built from the same proteins in the same cells, differing only in statistical form. \textbf{The winner swaps}: the mean incumbent leads under the mean-shaped oracle ($+0.242$ against $+0.146$), the distributional score under the distribution-shaped one ($+0.529$ against $+0.334$), in each of the three immune conditions. Hatched, the magnitude-matching control, which compares no distributions: it climbs from $+0.125$ to $+0.352$ as the oracle turns distributional, yet energy still leads it by $+0.177$ there against $+0.022$ under the mean-shaped oracle.
\textbf{Row 2}, the two gates on ZhaoSims2021, ten glioblastoma patients (Methods). \textbf{d}, Gate 2 posed as the same drug-response question in both settings (36 splits, 4 patients). Pooling drugs into classes holds the constructed ceiling at 0.700 (grey); one drug against one drug, the same cells give 0.879. All three arms are measured at the same granularity and under the same representation, so they are like for like; the 0.692 quoted for this construct elsewhere is the earlier in-fold run (Fig.~\ref{fig:6}e). In a real tumour the information is present (ceiling \textbf{0.923}) and unsupervised clustering does not reach it (\textbf{0.777}; median paired gap \textbf{$+$0.117}, 95\% CI $[+0.115, +0.135]$ by patient-level bootstrap, against $+0.007$ constructed; $+0.122$ against $-0.002$ with no method selection anywhere; $+0.123$ $[+0.098, +0.138]$ without the dominant patient). \textbf{Gate 2 in real tissue is an algorithmic bottleneck, not an informational one.} 30 of the 36 splits come from one patient and 39.9\% of cells fail the assignment margin, so both numbers are optimistic. \textbf{e}, Gate 1: induced response cosine, median \textbf{0.566} over 17 patient-drug pairs, against \textbf{0.014--0.044} in our constructed mixtures. \textbf{f}, The premise: over 18 within-patient drug pairs, mean-signature similarity ranks malignant-compartment response similarity at Spearman $\rho = \mathbf{\PREMISEOVERLAP}$, and the disjoint form of the same question, ranking the malignant compartment by the \emph{myeloid} compartment rather than by a mean that contains 43\% of the malignant cells, gives $\mathbf{\PREMISEDISJOINT}$. \emph{No $p$-value is given}, since 15 of the 18 pairs come from one patient. Both gates open here and the distributional advantage still does not appear: the two gates are necessary, not sufficient, and the condition that fails is the third one, decision relevance (Results).}
\label{fig:5}
\end{figure}

\subsection*{The information condition is a two-gate criterion}

The preceding results pose a mechanistic question: why is the distributional signal large
on observed populations yet absent under independent metrics? The answer is a
\textit{candidate-response information condition} with two distinct gates. A
distributional score can only exploit subpopulation structure whose subpopulations (i) actually
respond \emph{differently} in the candidate populations supplied to the retriever, and (ii) can
be told apart from those populations at all. Both gates are closed \emph{in the constructed
cell-line data this section uses}, which is why the observed-population result is an upper bound
rather than an operating point there. That qualifier is load-bearing and we did not originally
have it: on natural patient tumours both gates turn out to be open, and the last Results section
is the one that corrected us.

\textbf{A third condition, which we name here rather than measure as we measure the first two.}
Gates 1 and 2 concern whether the information a distributional score needs is present and
reachable. Neither asks whether it \emph{matters}. We therefore name a \textbf{Gate 3, decision
relevance}: the differential response must change \emph{which candidate is preferred} under the
outcome the decision is judged on, over and above what the mean signature already ranks. Its
status in this paper is different from the other two and we mark that difference wherever it
appears. Gates 1 and 2 are measured directly, on labelled mixtures and on patient tumours. Gate 3
is inferred, from three lines of evidence that were collected to answer other questions: in a real
tumour the response of one compartment already ranks the response of another at Spearman
$\rho = \PREMISEDISJOINT$, a value dominated by one patient (Fig.~\ref{fig:5}f); against the
independent functional oracle a scalar comparing no distributions reproduces most of the
distributional score's margin (Fig.~\ref{fig:4}h,i); and no true-divergence stratum shows a
positive mechanism-recovery gain (Fig.~\ref{fig:6}g). We propose Gate 3 as the condition that
explains why opening the first two is not enough, and we do not claim to have tested it as such.

\textbf{Gate 1, differential response.} We test what generated candidate populations actually
contain, using predict-then-rank experiments in which candidate responses are produced by a
perturbation predictor rather than observed (Fig.~\ref{fig:6}). Across three predictor families
the distributional advantage over mean retrieval is small and \emph{not consistent in sign}: the
nDCG@10 delta is $-$0.016 (average-effect), $-$0.071 (nearest-neighbour) and \textbf{$+$0.027}
(CPA-style linear-latent, a surrogate for the latent-arithmetic family; we do not benchmark a
trained scGen). The distributional score also wins on Hit@1 for two of the three. We therefore do not claim the
advantage is uniformly absent under prediction, only that it is small, unstable, and does not
survive the independent metrics of the previous sections.

The structural explanation is not the intuitive one. Predicted candidate populations do
\emph{not} collapse subpopulation structure. Synthesizing a predicted population by
applying the predicted effect to the query context's real control cells, which is what the
latent-arithmetic models in this class actually do, the subpopulation-variance ratio of
predicted populations is 0.142 against 0.138 for the real blended
candidate population that a scorer ranks; the latent model even sharpens it, to 0.462
(Fig.~\ref{fig:6}c). Structure survives.

What does not survive is \emph{divergence}, though the way we first measured it was doubly wrong
and the corrected version says something sharper. Measuring the induced response cosine between the
two subpopulations of a candidate population, $\cos(d_{\mathrm{maj}}, d_{\mathrm{min}})$ with each
subpopulation referred to its own matched control, real candidate populations respond
near-orthogonally (\textbf{0.014}) while predicted ones respond in largely aligned directions
(0.186 average-effect, 0.239 linear-latent, 0.366 nearest-neighbour). Those predicted
values do not mean what we took them to mean, for two reasons given below: the populations are
cross-context blends, and for a latent model the quantity is contaminated by the generator's
reconstruction error. We retain them because the structure-preservation conclusion above does not
depend on them, and we do not build any claim on them.

The limit of this behaviour is algebraic, not empirical, and the algebra is what carries the
argument. Call a predictor additive if it adds a \emph{single} delta vector to
\emph{every} cell of the query context, whether in gene space or through an invertible latent
map. Its two subpopulations then have \emph{identical} response deltas, $d_{\mathrm{maj}} =
d_{\mathrm{min}} = \delta$, so their induced cosine is exactly 1 and the divergence a
distributional score could exploit is exactly zero by construction. No amount of
training changes this. Additivity is thus a \emph{sufficient} condition for Gate 1 to be shut,
and it is a property of a model class, not a claim about any particular software.

Which deployed predictors fall in that class is a separate, empirical question, and the
within-context measurement can now answer part of it. Every number below is the model's \emph{true}
induced divergence, taken against its own baseline, and every model is first required to pass a
learning check: the cosine between its predicted mean delta and the true one. A model that
has not learned the drug effect cannot be said to open or shut any gate, and its divergence is noise
structure. We enforce that rather than assume it.

Real treated cells, split by the same control-derived cell state, give the target: their two
subpopulations respond at a cosine of 0.205, so within-context subpopulations in real data
do respond very differently. Against that, the average-effect predictor (learning check 0.502) and
the linear-latent model (0.343) both return 1.000000 exactly, reproducing none of it, as the
algebra requires. The published scGen variational autoencoder (learning check 0.420), which
this paper previously declined to say anything about, returns 0.972: its nonlinear decoder
buys a small departure from strict additivity, and it reproduces roughly 3\% of the divergence real
cells show. \emph{A deployed latent-arithmetic model is therefore near-additive in practice, not only
in principle}, and Gate 1 is essentially shut for it.

\textbf{A genuinely non-additive predictor does open Gate 1, but only a crack.} A predictor whose
displacement depends on the cell it is applied to should, if the two-gate account is right, show a
divergence below the additive ceiling of 1. We tested an entropic optimal-transport map, the
non-parametric core of the neural-OT perturbation models\cite{CellOT,Departures}, run in a latent
space as every deployed neural-OT model is. It passes its learning check ($0.352$, above the $0.30$
bar) and its \emph{true} induced divergence is $\mathbf{0.900 \pm 0.101}$, significantly below the
additive ceiling, so a non-additive predictor genuinely gives different subpopulations different
responses. But real treated cells sit at $0.205$, so the map closes only about an eighth of the
distance from the ceiling to what biology shows, and its learning is weak enough that part even of
that crack is prediction noise.

Two deployed models meet the same wall from opposite sides. The published CPA reconstructs cells
acceptably ($R^2 \approx 0.39$) but does not learn the drug-specific effect (predicted-versus-true
mean-delta correlation $0.08$), so it sits below the learning bar and its divergence is not
reportable; neither five times more cells per drug nor conditioning its perturbation embedding on
2048-bit Morgan fingerprints in place of a learned vector rescues it (learning $0.08$ and $0.03$),
which locates the failure in the adversarial objective rather than the drug representation.
CellFlow\cite{CellFlow}, an optimal-transport flow-matching model conditioned on the same
fingerprints, does learn within context (mean-delta cosine $0.24$ over twenty drugs, $0.81$ for the
HDAC inhibitor AR-42) and opens Gate 1 (induced divergence $0.85$). But the clean test requires
prediction into a held-out context, and there its learning collapses to $0.03$ at a matched budget.
That collapse is not an artefact of a mean-based check: re-scored with the gene-resolved and
distributional metrics used by community virtual-cell evaluations\cite{VirtualCellChallenge}, the held-out-context predictions match the true differential-expression
direction only at chance ($0.52$ against $0.70$ within context), recover $7\%$ of the true
differential genes, and sit farther from the true treated population than the unperturbed control
does (energy ratio $2.1$ against $0.98$). It is instead mechanistic: per drug, leave-one-context-out
learning tracks that drug's cross-context response similarity (Pearson $r = 0.85$), so a predictor
conditioned on chemistry alone recovers only the context-shared part of a response, which on this
dataset is small (median cross-context cosine $0.04$) and near zero for exactly the
response-divergent drugs the retrieval task selects.

\textbf{Whether opening Gate 1 buys retrieval gain therefore stays open, for an instructive reason.}
Re-running the cross-line predict-then-rank with the OT map as the candidate-response source, its
predictions are too poor to retrieve with at all: Hit@1 is $0.05$ under \emph{every} rule, mean
cosine included, against $0.88$ for the additive average-effect predictor. \textbf{We lack a
predictor that both predicts accurately across contexts and opens Gate 1}, which is exactly what a
clean test of the second half of the two-gate account needs. That we could not build one is
consistent with the rest of this paper, in which the accurate part of every predictor is its mean
and the part that departs from the mean is noise, but it is a limitation rather than a falsification
and we present it as such.

\textbf{One measured value is an artefact of the generator, not of the biology.} An additive
predictor is pinned by the algebra to return exactly $1.000000$, and asserting that value is what
exposed two errors in our own first design: the candidate populations were cross-context blends,
and a latent model's response must be referred to the right baseline (Methods). The second of those
is a result rather than a correction. A latent model does not add a delta to a cell, it
\emph{reconstructs} the cell and adds a delta to the reconstruction, returning $P(C) + W\delta_z$.
Referred to its own decoded baseline the reconstruction term cancels and the induced cosine is
exactly $1.000000$, as the theorem requires. Referred to the \emph{real} control cells, which is
what a retrieval scorer necessarily does, it carries each subpopulation's reconstruction bias,
which is cell-state dependent and does not cancel, and the model then appears to induce a great
deal of divergence while inducing exactly none. \textbf{A distributional score run on such a
candidate population is confidently scoring an autoencoder's artefacts}, and that includes the
structure diagnostics reported above whenever they are taken against real controls.

\textbf{Gate 2, recoverability.} Even where differential structure exists, a decision layer must be
able to resolve it. We test this on a controlled mixture with ground-truth labels (K562 at
10\,$\mu$M, HDAC-class versus JAK-class cells), where the bimodal structure exists by construction.
The question that matters is not whether clustering succeeds but \emph{why it fails}, and those are
two very different papers: if the information is absent from the representation no method can help,
whereas if it is present and only the unsupervised assignment step cannot reach it, the bottleneck
is an algorithm and is solvable. Distinguishing them requires an upper bound, which prior work in
this area, ours included, has not reported. We therefore score every method in one unit,
best-permutation accuracy against the true labels, with label matching given to each clusterer for
free (Fig.~\ref{fig:6}e).

Unsupervised recovery is poor, the best of four methods reaching $0.674$, but the informative number
is the \textbf{supervised ceiling}, which reaches only \textbf{0.692} (AUC 0.752); a nonlinear
learner does not beat a linear one, so the boundary is not a hidden nonlinearity our probe was too
weak to find. \textbf{The gap between the ceiling and the best unsupervised method is
$\mathbf{+0.018}$ accuracy}, so \emph{within this benchmark} clustering already extracts essentially
everything the representation holds and the limit is informational rather than algorithmic. Two
later results retract almost everything one might conclude from that sentence, and we flag them here
rather than let a reader carry it forward. The gap is an \emph{underestimate}, because the design
gives the clusterers the easier transductive representation and the probe the harder in-fold one
(we had stated that bias backwards; Methods). And the $0.692$ ceiling is largely an artefact of
\emph{pooling} several HDAC and several JAK compounds into two classes: split one drug against one
drug, the same cells give $0.879$ (Fig.~\ref{fig:5}d). The informational reading survives neither
correction, and it does not survive contact with tissue. It is nonetheless a measurement rather
than an insensitive probe: raising the separation between the two states artificially makes the
ceiling pull away from clustering exactly as it should (Fig.~\ref{fig:6}f), and across 529
(cell line, drug) pairs in SciPlex3 the median silhouette of a two-way split is $0.034$ (Extended
Data Fig.~3).

\textbf{But this benchmark is one we built, and the next section shows that it misled us.} An
HDAC-treated and a JAK-treated K562 cell are a mixture we assembled; they are not a subpopulation
structure that occurs in a patient. Read on its own, the number above would license the conclusion
that real subpopulations sit at an information ceiling no method can exceed, and we drew exactly
that conclusion before testing it. It is wrong.

\textbf{What Gate 2 does and does not constrain.} Energy distance, MMD and sliced-Wasserstein are
\emph{cluster-free}: they compare whole populations and never form a partition, so Gate 2 is not a
claim that they require one. The constraint is subtler, and it closes the loop with the
functional-oracle result. A distributional distance sums every source of population difference into
a single number, so where the subpopulation signal sits at a ceiling as low as the one measured
here, that signal cannot be what dominates the distance. Something else must, and we have measured
what: the candidate's own response magnitude, which tracks the energy distance at
$\rho = +0.791$. Gate 2 does not stop a cluster-free score from being computed; it determines
whether the number that score returns is about subpopulation structure or about something else.

The two gates compound, and their consequences differ by score type. Note what Gate 1 is
\emph{not}: it is not structure preservation. Every predictor we tested preserves structure and
every one still fails, so a generator validated on a structure-preservation metric would pass a
test that does not bear on the question. Coverage-type scores, which need the partition explicitly,
are blocked wherever Gate 2 is shut, which in this study is the constructed mixtures and not the
patient tumours of the next section. Global scores are never blocked, but are left measuring
whatever channel is largest, which is not the one their advocates intend.

Two further tests, both consistent with the criterion and neither rescuing it, are given in
Supplementary Note 2: coverage performance rises with query identifiability but only weakly and
not monotonically, and when queries are stratified by the Gate-2 variable itself the
distributional advantage is flat.

% FIGURE 5
\begin{figure}[!htbp]
\centering
\includegraphics[width=\textwidth]{figures/fig6.pdf}
\caption{\textbf{When can distributional information matter? Two measured gates, and a third condition they do not cover.} Candidate populations must show differential response (Gate 1) and that structure must be identifiable (Gate 2); neither asks whether the difference changes the decision (Gate 3). \textbf{a}, The predict-then-rank design and the conditions a candidate population must satisfy: a held-out query is scored against candidate populations that are either observed or produced by a perturbation predictor. Gates 1 and 2 are measured in panels \textbf{b}--\textbf{f}. \textbf{Gate 3 is drawn dashed because it is proposed here rather than tested as the other two are}; its evidence is Fig.~\ref{fig:5}f and Fig.~\ref{fig:4}h,i, and it is what panel \textbf{g} is consistent with. \textbf{b}, Predictor-specific retrieval gain (nDCG@10, distributional minus mean). The gain is small and \emph{not} consistent in sign: $-0.071$ (nearest-neighbour), $-0.016$ (average-effect), $+0.027$ (latent-linear); one task, so no dispersion is available. \textbf{c}, Gate 1. Predicted candidate populations \emph{retain} baseline subpopulation structure (variance ratio 0.142 versus 0.138 for real ones, left axis) but do not reproduce \emph{divergence}: the induced response cosine is 0.19--0.37 for predicted populations against 0.014 for real ones (right axis), where an additive predictor is pinned at exactly 1 by algebra. \textbf{The predicted values here do not measure that ceiling and we do not interpret them}, because these populations are cross-context blends and a latent model's response against real controls also carries its reconstruction bias (Methods). \textbf{d}, Structure diagnostics, predicted relative to real. \textbf{e}, Gate 2 on a known bimodal mixture (K562, HDAC-class versus JAK-class cells), every method scored in one unit, best-permutation accuracy against the true labels. Unsupervised recovery is poor (best 0.674) and the \emph{supervised ceiling} reaches only 0.692 (AUC 0.752), a gap of \textbf{0.018}, so \emph{within this benchmark} clustering already extracts what the representation holds. That reading is superseded: the ceiling is largely an artefact of \emph{pooling} drugs into classes, and in a patient's tumour the same question gives 0.923 against 0.777 unsupervised, an algorithmic limit rather than an informational one (Fig.~\ref{fig:5}d). \textbf{f}, The control that licenses reading panel \textbf{e}: raising the separation between the two states makes the ceiling pull away from clustering as it should, so the absence of a gap at real separation is a measurement rather than an insensitive probe. \textbf{g}, The conditional advantage a working information condition would produce is absent. Distributional-minus-mean MoA-nDCG gain by \emph{true} response-divergence quartile (bars are means, error bars s.e.m.; Benjamini--Hochberg corrected). No stratum is positive and the lowest-divergence quartile is significantly \emph{negative} (mean $-0.082$, $q = 5.8\times10^{-5}$, $n = 166$). Where the theory predicts the largest gain there is none.}
\label{fig:6}
\end{figure}

\subsection*{HIR-Bench separates observable from circular predictors}

To study these regimes under controlled conditions we construct HIR-Bench (Heterogeneous
Inverse Retrieval Benchmark), a synthetic, method-neutral generator with a specified
latent welfare oracle (Fig.~\ref{fig:5}). Each task instantiates one of three information
conditions, observed, predicted\_mean (candidates collapsed to their
means), and predicted\_structure (candidates with partially preserved structure),
and the generator admits an analytical flip boundary $\alpha^* = B/(A+B)$ separating the
regime where distributional retrieval can help from where it cannot (Fig.~\ref{fig:5}a).

The grid comprises 13,440 benchmark \emph{instances} (672 parameter configurations
$\times$ 20 seeds; these are synthetic retrieval instances, not biological cells). The
sanity checks hold (no-conflict flip rate 0.0000; median boundary margin 0.27; Extended
Data Fig.~6d), and collapsing candidates to their means inflates decision regret most
severely against the worst-welfare objective, the one that depends on the minority
subpopulation (0.29 observed versus 0.96 predicted-mean; Extended Data Fig.~6c).

\textbf{Is retrieval failure predictable, and from what?} The second half of that question
matters more than the first, and a benchmark with a latent oracle is the only place it can
be answered cleanly. Features derived from the oracle are available to the benchmark's
authors but never to a retrieval method at query time; features of the observed
populations are available to both. Fitting the same model on each feature set separates a
usable diagnostic from a circular one, and the two answers are opposite.

We therefore run a $2\times2$ design: two feature sets crossed with two cross-validation
groupings. The \emph{observable} features (subpopulation-variance ratio, isotropy and response
diversity of the query and candidate populations) are computable at query time. The
\emph{oracle-derived} features (top-$k$ disagreement, Kendall conflict) are each a function of
the benchmark's latent utility matrix, as is the label. The groupings differ in their unit: the
label is a deterministic step function of $(\alpha, \mathrm{conflict})$, so the 13,440 instances
contain only 28 independent parameter cells, whereas the natural leave-one-instance-out
scheme treats 672 as independent and leaves roughly twelve near-identical copies of each held-out
instance in training.

All four cells are shown in Fig.~\ref{fig:5}b. Read across the top row, retrieval failure
\emph{is} moderately predictable from what a method can see: AUC \textbf{0.788} against a
majority-class rate of 0.643, the query's response diversity carrying the highest weight. Read
across the bottom row, the oracle-derived features are no better than chance (0.400;
with 28 independent cells this is not distinguishable from 0.5). They predict one function of the
latent utility matrix from another, which measures circularity, not predictability.

The two columns separate the two failure modes cleanly, and the interaction is the point.
Pseudo-replication inflates the circular feature set by $+$0.240 (0.400 to 0.640)
but the honest feature set by only $+$0.048 (0.788 to 0.836). An instance-level
holdout preferentially rescues features that carry no generalizing signal, because
memorizing near-duplicates is all they have; features with real signal barely benefit.
A benchmark whose cross-validation unit is finer than its label-determining axes will
therefore manufacture apparent predictability precisely where there is none, and the effect
is largest for exactly the features an author is most tempted to use. Neither effect is
specific to this benchmark.

Two caveats bound the 0.788. The
effective sample size is 28 parameter cells, so a confidence interval computed on
$n = 13{,}440$ would be meaningless, and by the same token the gap between 0.400 and 0.5 is
not itself interpretable. And the label is a deterministic step function of those
parameters, which makes every held-out fold single-class, so no per-fold AUC distribution
exists and the pooled out-of-fold AUC is the only statistic available. ``Moderately
predictable'' is as strong a claim as this supports.

Two limitations are essential. HIR-Bench is intentionally constructive: built from a specified
latent welfare model, its phase boundary is expected by design, and the question that
matters is whether real data project onto the same regimes. They do not: the boundary fails to
transfer, scoring 61.1\% against a 95.4\% majority-class baseline (Methods). And one apparently
supportive number is not evidence: the gate ``correctly'' flags mean-collapsed predicted responses
as no-DART in 100\% of the transfer test, but that rate is a tautology, since the rule
returns no-DART unconditionally for that information condition (Supplementary Note 1). HIR-Bench is
a controlled framework for exposing failure modes, not a real-world diagnostic.

\subsection*{In real tissue Gate 2 is an algorithmic limit}

Every real-data positive in this study so far lives in cell-line mixtures that \emph{we}
constructed. A minority subpopulation assembled by pipetting A549 together with K562 is not
evidence about the heterogeneity a tumour actually has, and a criterion validated only on such a
mixture is a criterion validated on our own pipetting. We therefore tested the two gates on
ZhaoSims2021, acute-slice culture and surgical biopsy from ten glioblastoma patients, in
which malignant glioma cells, tumour-associated myeloid cells and oligodendrocytes co-exist inside
one patient's tumour. Nobody mixed anything. Compartments are assigned by validated marker score
and cells without a confident call are dropped rather than reassigned; an earlier attempt that took
a bare argmax over five marker sets fabricated a 22\% ``T-cell'' compartment whose cells did not
express CD3 (Methods). This is \emph{not} a retrieval benchmark and we do not use it as one: the
panel is six drugs and only one patient received more than two.

\textbf{Gate 2 must be measured on the same construct in both settings, and our first attempt was
not} (Fig.~\ref{fig:5}d). The constructed Gate-2 benchmark separates HDAC-treated from JAK-treated
K562 cells: one lineage, cells differing only in which drug they saw, a \emph{drug-response}
partition. Our first natural test separated malignant from myeloid \emph{control} cells: two
lineages, a \emph{cell-type} partition, and one that is trivially easy because a glioma cell and a
macrophage differ in thousands of genes. It reached a ceiling of 0.964, and comparing that to
0.692 compared two different questions. Gate 2 is defined on the partition a retrieval score must
resolve inside a candidate population, which is a drug-response partition. We therefore re-ran it
as one: cells treated with drug A against cells treated with drug B, \emph{within one compartment,
within one patient}, which is the K562 construct transplanted into tissue nobody assembled
(36 drug-pair splits; Methods). We additionally report the ceiling under the representation the
clusterers actually receive, because the in-fold representation used earlier deflates the ceiling
and so \emph{shrinks} the very gap the conclusion rests on.

\textbf{Two things follow, and the second one reverses us.} First, the constructed ceiling was
substantially an artefact of \emph{pooling drugs within a class}: run at this comparison's own
granularity and representation the pooled construct gives 0.700, against the 0.692 the earlier
in-fold run of the same construct reported. Splitting the same
K562 cells one drug against one drug, as the natural test necessarily does, lifts the ceiling to
0.879 with unsupervised clustering at 0.837. Much of what we read as an
information limit was the within-class heterogeneity of several distinct compounds we had merged.

Second, and this is the finding: in a real tumour the information is plainly \emph{present}, and
unsupervised clustering \emph{cannot reach it}. The supervised ceiling is 0.923 while the
best unsupervised method reaches only 0.777, a gap of $+$0.117 against $+0.007$
on the like-for-like constructed split. Centroid-based methods collapse ($k$-means 0.638, Gaussian
mixtures 0.620) where the graph-based method partly survives (Leiden 0.780), which is the signature
of a dominant variance direction that is not the drug-response axis: inside a real compartment,
cell-state heterogeneity outweighs the drug response and unsupervised partitioning locks onto it.

\textbf{We therefore withdraw the information-limit claim, and we withdraw the first withdrawal
too.} The original conclusion, that real subpopulations sit at an information ceiling no method can
exceed, is wrong. But so is the reading we replaced it with, that real subpopulations are simply
easy to resolve; that came from the cell-type comparison and does not survive being posed as a
drug-response question. In real tissue Gate 2 is an \emph{algorithmic} bottleneck, not an
informational one: the partition is there to be found, off-the-shelf clustering does not find it,
and the gap is therefore something a better method could close. That is the one place in this study
where the news for distributional retrieval is constructive rather than cautionary, and it is the
opposite of what we concluded from our own benchmark twice.

\textbf{We attacked that gap three ways before believing it, because it is the only constructive
claim in this paper.} It is a difference of two \emph{maxima}, so we recomputed every fixed
(probe, clusterer) pair with no selection anywhere; the 36 splits are not independent, since 30
come from one patient, so we bootstrap over \emph{patients} rather than splits, the
pseudo-replication this paper criticises elsewhere; and we dropped the dominant patient entirely.
The gap survives all three, and dropping that patient makes it slightly \emph{larger}, not smaller
(values in Fig.~\ref{fig:5}d; none of 5,000 bootstrap draws lies at or below zero).

\textbf{Three limits nonetheless bound it.} A cluster bootstrap with only \emph{four} clusters
under-covers, so the interval is a floor on the uncertainty rather than an estimate of it.
\emph{Cells were dropped}: 39.9\% fail the compartment-assignment margin and are excluded rather
than reassigned (Methods, Supplementary Table 3), and since the discarded cells are by construction
the hardest to place, both the ceiling and the clustering accuracy are optimistic. \emph{The drug
panels differ}: the constructed split contrasts HDAC with JAK inhibitors and the natural one six
unrelated compounds, so response magnitude is not matched across the arms. We report the comparison
because it is the closest to like-for-like these data permit, and claim no more precision for it
than four patients and one drug panel can support.

\textbf{Gate 1 opens too, and the same benchmark misled us in the opposite direction}
(Fig.~\ref{fig:5}e). Naturally co-existing compartments do respond in different directions to the
same drug, but far less so than our mixtures implied: the induced response cosine has a median of
0.566 across 17 patient-drug pairs, where our constructed mixtures sit at
0.014--0.044, essentially orthogonal. One drug pushes two unrelated cell lines in unrelated
directions; malignant cells and tumour-associated macrophages in one patient share most of their
response. The constructed benchmark overstates the divergence a distributional score can
exploit, by roughly twofold on a divergence scale ($1-\cos$: 0.99 against 0.43), while misplacing
the identifiability limit. (The ratio of the cosines themselves, 40-fold, is not a meaningful
measure of divergence when one cosine is near zero.) Both distortions
flatter this paper's premise, and neither was intentional.

\textbf{Both gates are therefore open in real tissue, and the two-gate account predicts that a
distributional advantage should follow. The mechanism that would produce one is missing, and its
absence is what Gate 3 names.} The
paper's motivating scenario, that two drugs can share a mean signature and do opposite things to a
subpopulation, can be tested directly on a real tumour. Across 18 within-patient drug pairs, the
similarity of two drugs' \emph{mean} signatures ranks the similarity of their
\emph{malignant-compartment} responses at Spearman $\rho = \mathbf{\PREMISEOVERLAP}$
(Fig.~\ref{fig:5}f).
Where the mean says two drugs are alike, the malignant compartment mostly agrees. Of the five drug
pairs most similar on the mean, two have a malignant response meaningfully less similar than their
mean suggests, so the residual is real; it is also small.

\textbf{That statistic has a mechanical component, and we report the corrected one alongside it.}
The mean signature is taken over all called cells and malignant glioma is 41,314 of the 96,225 of
them, so the two quantities being correlated share 43\% of their cells and part of the agreement
is arithmetic rather than biology. Recomputed on compartments that share \emph{no} cells, ranking
the malignant-compartment response by the \emph{myeloid}-compartment response instead of by the
mean, the correlation is $\rho = \mathbf{\PREMISEDISJOINT}$ (\PREMISEOVERLAP\ to
\PREMISEDISJOINT{}; on PW030 alone, $+0.821$ to $+0.804$). The finding survives the correction and
the number moves, which is the honest way round. \emph{We treat the disjoint value as the
load-bearing one}, because it is the one that answers the question actually posed: whether what
the bulk of the tumour does already tells you what the compartment of interest does. A paper
arguing that a scoring rule and its evaluator must not share their object cannot exempt its own
premise statistic from the same rule.

This is as close to a direct reading of
Gate 3 as these data allow: where the mean already ranks what the subpopulation does, the
differential response has little left to change about the decision. The two limits below are what
stop it from being a measurement of Gate 3 rather than an indication of one.

\textbf{Two limits on that number, and they are severe.} It is close to a \emph{single-patient
result}: 15 of the 18 pairs come from PW030, the only patient given more than two compounds, and
reuse the same six drugs, while three patients contribute one pair each. A $p$-value computed as
though the 18 pairs were independent would be anticonservative by a wide margin, so we do not quote
one; the honest description is a strong association in one patient, consistent in sign with three
single-pair observations elsewhere. Testing the premise properly needs a cohort in which many
patients receive many drugs, which this dataset is not. And a rank correlation between two
similarity scores is neither a fraction of variance nor a statement that the response vectors are
collinear, and we read it as neither.

\textbf{We did not run retrieval here, so we report no retrieval result.} ZhaoSims2021 supports six
compounds with only one patient receiving more than two, and a leave-one-drug-out Hit@1 computed on
that would be an anecdote wearing a benchmark's clothes. What this section shows is that the
\emph{mechanism} a distributional advantage would run through, mean-similar drugs diverging inside
a subpopulation, is largely absent in real tissue: an inference about why an advantage should not
be expected, not a measurement of its absence.

\textbf{The two gates are necessary but not sufficient, and this is how we found out.} Opening them
does not deliver a distributional advantage, because in real tissue the compartment-specific
response is mostly ranked by the mean already. The independent functional oracle reaches a
compatible conclusion by an entirely different route: energy retrieval beats the mean incumbent
there, and a scalar magnitude control then reproduces almost the whole of that margin. The two
lines of evidence share no data, no metric and no code path, and they do not share a number
either, so we do not claim they converge on one. Both are consistent with Gate 3 being the binding
one here, and neither was designed to test it. What they agree on is a direction: distributional
information in this domain is real, positive, and much smaller than objective-aligned evaluation
makes it appear.

% FIGURE 6

%% =====================================================================
%%  DISCUSSION
%% =====================================================================
\section*{Discussion}

The main result of this work is not that DART is better or worse, but that \emph{how much} better
it is depends almost entirely on who is asked. Under a metric of its own distributional kind the
advantage is large. Under an independent, semantically matched functional oracle it is real but
small, and almost all of what remains is a response-magnitude channel that a scalar reproduces
without comparing any distributions at all. Under a task-proximal mechanism criterion it is
absent. The same rankings, the same scorer, the same queries: only the judge changes. DART is
valuable here as a probe, an instrument that can be scored by its own objective and by
more independent criteria on the same tasks, rather than as a therapeutic recommender. The paper's
positive anchors survive intact: the unification of mean and distributional retrieval is an
algebraic fact, and there is genuine distributional signal in response-available populations. What
does not survive is the size of the gain that objective-aligned evaluation reports for it.

We also asked, as a direction for future work rather than a load-bearing result, whether the
response divergence a distributional score measures has any biological correlate. In the Frangieh
Perturb-CITE-seq data, cells whose response direction diverges most between contexts are weakly
enriched for resistance-associated programs. Without a drug timecourse, a survival readout or a
post-treatment clone measurement this is hypothesis-generating only, we make no claim that
DART predicts resistance, and the analysis is reported in Supplementary Note 3 and Extended Data
Fig.~4 rather than here.

\subsection*{Why the negative results are central}

The predict-then-rank nulls and the independent-metric results are the evidence that supports this
paper's core claim, not a shortfall against it. The objective-aligned positive on its own would
license the conclusion that distributional retrieval substantially improves ranking. The
task-proximal metrics do not support that, and the independent functional oracle supports a far
weaker version of it: the distributional score does beat its incumbent there, but a scalar that
performs no retrieval reproduces nearly the whole margin. The positive and the null together,
labeled by metric class, are what tell a practitioner which conclusions a given metric licenses
and which it does not.

\subsection*{Implications for virtual-cell and perturbation-prediction models}

The two-gate result has a direct implication for the fast-growing class of perturbation predictors
and virtual-cell foundation models. In this framework such models are
candidate-response sources, and whether a distributional decision layer can add anything
on top of them depends entirely on whether they supply \emph{differential} response. The
load-bearing statement here is algebraic and needs no benchmark to establish: a predictor that
adds one delta vector to every cell induces exactly zero response divergence, so additivity is
\emph{sufficient} to make a distributional decision layer worthless on top of it. That is a
property of a model class, not a deficiency of any implementation, and no amount of training
removes it.

What we do not claim is that today's predictors are, as a field, additive. We can say it of
one deployed model: the published scGen variational autoencoder passes its learning check and sits
at a true induced divergence of 0.972 against the additive ceiling of 1.000, reproducing about 3\%
of the divergence real subpopulations show. That is a fact about scGen, not about the field. The
in-repository linear-latent model is additive by construction, so its value demonstrates the algebra
rather than surveying anything.

Non-additive, cell-state-conditional predictors exist, including optimal-transport, flow and
Schr\"odinger-bridge formulations\cite{CellOT,Departures}, and a latent-space optimal-transport map
does open Gate 1, though only a crack. What we have not done is the second half of the test.
A predictor that genuinely produces state-specific responses should, if the two-gate framework is
right, restore a distributional advantage that additive predictors cannot, and one that opens Gate 1
yet yields no such advantage would falsify the framework. That test is out of reach for the
predictors we could run: the optimal-transport map's cross-line predictions are too weak to retrieve
with, and the flow model this test most plausibly calls for, CellFlow\cite{CellFlow}, learns and
opens Gate 1 within a context but does not carry that learning to a held-out cell line, because per
drug the transfer is bounded by the drug's cross-context response similarity and the retrieval task
selects exactly the drugs for which that similarity is near zero (Results). The decisive test
therefore still requires a predictor that both predicts accurately \emph{across contexts} and opens
Gate 1, most plausibly a flow or Schr\"odinger-bridge model trained to transfer; we name this
as the test that would falsify us, and record that the flow model we trained does not yet complete
it.

The methodological point stands independently of that test, and it is sharper than we first
realised. Differential response, not mean-level accuracy and not even structure preservation, is the
property that determines whether distributional retrieval is applicable at all. Structure-preservation
metrics alone would have scored every predictor we tested as adequate. Worse, they would have
scored the generator's own reconstruction error as biological structure. An additive latent model
induces exactly zero true response divergence, yet measured against real control cells, as a scorer
must, it appears to induce a great deal, because the autoencoder reconstructs the two subpopulations
with different biases. A distributional decision layer stacked on such a model would be scoring the
autoencoder. Anyone reporting that a generative perturbation model ``preserves subpopulation
structure'' or ``induces divergence'' should therefore state which baseline the claim is measured
against, and show that it survives being measured against the model's own output. This
is consistent with independent benchmarking, in which foundation models and deep-learning
predictors did not outperform deliberately simple baselines for perturbation
prediction\cite{ref13,ref14}: objective-aligned reconstruction accuracy need not translate into an
advantage over mean or linear baselines.

\subsection*{When distributional information can matter}

The question a practitioner actually asks is not whether distribution-aware retrieval is
defensible in principle but when it is worth reaching for, and the results here answer it as three
conditions that must hold together:

\begin{quote}
\emph{heterogeneous responses exist} (Gate 1, differential response) $+$
\emph{those responses are recoverable} (Gate 2, recoverability) $+$
\emph{they change which candidate is preferred} (Gate 3, decision relevance)
\end{quote}

\noindent Each is necessary and none is sufficient, which is the practical content of the result:
we found a setting, patient glioblastoma, in which the first two hold and no advantage follows.
The three are not on equal evidential footing and we do not present them as though they were.
Gates 1 and 2 are measured here, on constructed mixtures and on patient tumours, and Gate 1 in
addition has an algebraic form that needs no benchmark: an additive predictor shuts it by
construction. Gate 3 is proposed rather than measured, and rests on evidence gathered to answer
other questions, chief among it a within-patient correlation dominated by a single patient
(Results). A study reporting a distributional gain can nonetheless check all three, and the third
is the cheapest of them: show that the ranking the distribution produces differs from the ranking
the mean produces on the outcome being claimed. Where it does not, the distribution is being
computed but is not deciding anything.

\subsection*{What independent validation would require}

The external functional readout we could reach, GDSC dose-response profiles matched to our drug
panel, constrains the question without settling it. It is a semi-real oracle: viability is
measured in a separate bulk assay, matched by compound name, under a different dose regime. It
establishes that transcriptional similarity recovers functional similarity better than a mean
signature does, and that almost all of that recovery is a magnitude channel. It cannot establish
therapeutic utility. Turning DART from an evaluation probe into a therapeutic method would require
readouts measured on the same cells the method is asked to treat: dose-response viability
\emph{in situ}, longitudinal treatment survival, functional perturbation validation,
post-treatment clone expansion, and combination-response validation. Until then, the defensible
claim is that distributional information is \textit{present} and \textit{modestly useful}, and far
smaller than objective-aligned evaluation makes it appear.

We urge four specific controls on anyone reporting a validation of a distributional retrieval
method, each of which caught an error of ours.

First, match the oracle to the task. A similarity retriever graded on absolute potency is
being asked a question it was never posed, and it will fail in a way that looks like a finding:
potency tracks response magnitude, a distributional distance also tracks response magnitude, and
the two confound. We ran that comparison, drew a conclusion from it, and had to withdraw the
conclusion; it now appears only as an audit.

Second, beat the magnitude scalar. Rank the candidates by how closely their response
magnitude matches the query's, compare no distributions at all, and show that your method beats it.
On the matched oracle ours does not do so significantly, and on the mismatched one it loses
outright. A distributional method that cannot clear this bar has not shown that the distribution is
what is doing the work.

Third, vary the shape of your external oracle, not just its identity. It is natural to
assume that once a criterion is external, independent and invisible to the method, it can be
trusted to arbitrate. It cannot. Two oracles built from the same proteins in the same cells, one
computed as a mean and one as a distribution, hand victory to opposite methods (Fig.~\ref{fig:5}c).
A field that validates distributional methods against distributional readouts and mean-based methods
against mean readouts will find both vindicated, and neither finding will be about biology. The
practical form of this control is cheap: compute your external criterion in at least two statistical
forms and report both. If your method only wins under the form that matches its own, say so.

Fourth, and least comfortable, do not let a benchmark you built decide a claim about
biology, and check that the two things you are comparing are the same question. Our constructed
cell-line mixtures simultaneously understated how identifiable real subpopulations are (ceiling
0.879 on a like-for-like split, against 0.923 in a patient's tumour, where our class-pooled
mixture had said 0.692) and overstated how divergently they respond (cosine 0.014 against 0.566).
Both distortions happened to favour the premise we were testing, and neither was visible from
inside the benchmark. Worse, our first attempt to check the benchmark against tissue compared a
drug-response partition in the mixture with a \emph{cell-type} partition in the tumour, which is a
far easier problem, and it produced a reassuring number (0.964) that meant nothing. A comparison
across two settings is only as good as the sameness of the question asked in each, and it is
alarmingly easy to hold the protocol fixed while silently changing the question. Any
information-condition criterion validated solely on synthetic or constructed heterogeneity should
be treated as a statement about that construction until it is checked against tissue, and the check
must ask tissue the same question.

\subsection*{Field-level evaluation implications}

The broader lesson is that the central distinction is not distributional versus mean retrieval, but
objective fidelity versus independent utility. Our field-level audit shows that
objective-aligned evaluation is the default across signature retrieval, expression
prediction, and graph-based inverse design: the headline metric is usually a restatement of
the method's own objective. This is legitimate for testing objective fidelity, and we do not
claim any prior method is wrong. But it means that strong benchmark numbers, on their own,
answer a narrower question than they are often read to answer.

Two corollaries follow. A benchmark whose cross-validation grouping is finer than its
label-determining axes will manufacture apparent predictability from nothing, as ours does for
oracle-derived features\cite{Leakage,PseudoRep}; and a distributional score's apparent link to therapeutic value can be a
response-magnitude confound, so any such claim should be required to beat a query-independent
magnitude baseline. Neither failure mode is specific to this work. This reading is corroborated by
recent metric-focused studies: the distributional metrics used here have documented failure modes,
with the Wasserstein distance degrading under variance scaling in high-dimensional expression space
and the energy distance overlooking disruptions in gene--gene dependencies\cite{ref16}, and metric
choice alone can flip model rankings\cite{ref17}. We cite the first against ourselves; it bounds the
interpretation of our own Class-A scores. The contribution is to make the
objective-fidelity/independent-utility distinction explicit and measurable on a single task set, a
template other inverse-retrieval methods could adopt.

\textbf{The same distinction may matter outside retrieval, and we raise it as a question rather
than as a result.} Objective--utility mismatch as we define it is a property of retrieval
evaluation, and it is the only place we have measured it. The ingredient it needs is not peculiar
to retrieval, however: a surrogate objective standing in for a functional outcome nobody measured.
Virtual-cell and perturbation-prediction models are scored on held-out reconstruction of the
quantity they were trained to reconstruct, and whether that tracks downstream decision quality is
a separate question that independent benchmarks have already begun to
raise\cite{ref13,ref14,VirtualCellChallenge}. Biomarker discovery selects features by an
association statistic and is frequently validated by the same statistic in a new cohort, which
tests transportability rather than clinical utility. Molecular design optimises a computed
property and is commonly evaluated by that property recomputed. We have tested none of these and
claim nothing about them. What transfers is the check, not the finding: report a result under at
least one criterion the method's own objective did not define, and say which class the headline
number belongs to.

\subsection*{Limitations and scope of claims}

A full statement of what this study does \emph{not} claim, seventeen items covering the scope of
every result reported here, is given as Supplementary Note 4. The load-bearing ones are these. No
gain in this study is measured on a clinical or same-cell endpoint. The Class-C positive is
conditional on the cell-line centering of its oracle and does not survive the uncorrected
construction, and under neither construction does the distributional score significantly beat a
scalar that compares no distributions. We do not claim that the predictors tested here represent
the field, nor that opening Gate 1 buys retrieval gain, since we could not build a predictor that
both transfers across contexts and opens it. And we do not report a count of distributionally
dominant real tasks, because the conventional threshold sits inside the noise band of the
difference distribution.

Three limits bound the scope of the whole study rather than any single result, and we state them
as such.

\textbf{First, the external oracle is imperfect, and its imperfection is not incidental to the
argument.} The Class-C readout is semi-real: viability measured in a separate bulk assay, joined
to the transcriptional data by compound identity rather than by shared cells, under a different
dose regime. A paper whose thesis is that the choice of criterion decides the winner cannot exempt
its own criterion, and we do not: the affirmative Class-C result holds under the cell-line-centered
oracle we argue is correct and reverses under the uncorrected one, and both are given in
Supplementary Table 4. The oracle-shape test makes the same point sharper, since two external
criteria built from the same proteins in the same cells hand victory to opposite methods.

\textbf{Second, this study audits retrieval, not prediction.} Every result here concerns how
candidates are \emph{ranked} and how that ranking is \emph{graded}. Perturbation predictors enter
only as candidate-response sources, and the one statement we make about a model class, that
additivity shuts Gate 1, is algebraic rather than a benchmark of any implementation. Nothing here
measures whether distributional \emph{prediction} objectives improve prediction, which is a
different question with its own literature\cite{CellOT,Departures,CellFlow}.

\textbf{Third, prospective validation remains outstanding, and it is what would settle the
question.} No endpoint in this study is measured on the cells a method is asked to treat.
Converting the conditional claim made here, that distributional information is present and modestly
useful, into a therapeutic one requires same-cell readouts: dose-response viability \emph{in situ},
longitudinal treatment survival, and post-treatment clone expansion. Until such a readout exists,
the affirmative content of this paper is bounded by the semi-real oracle that produced it.

%% =====================================================================
%%  METHODS
%% =====================================================================
\section*{Methods}

\subsection*{Data sources and preprocessing}

Three single-cell datasets carry every result reported here. SciPlex3\cite{dataSciPlex3}: 276,325 cells $\times$ 2,000 genes across three contexts (A549, K562, MCF7) and 188
drugs, all present in all three contexts. Frangieh melanoma Perturb-CITE-seq\cite{dataFrangieh}: 218,023 cells $\times$ 2,000 highly variable genes across Control,
IFN$\gamma$, and Co-culture conditions with 239 CRISPR knockouts. CD34$+$
(GSE306429)\cite{dataCD34}: 33,984 cells $\times$ 2,000 genes, 36 drugs, used as an honest negative
control. A fourth resource, the LINCS closed-loop / PDGrapher benchmark (41,070 rows
across A549, MCF7 and PC3, 1,369 drugs), was prepared for a graph-based inverse-design comparison
that we do not report; \emph{no number in this paper is derived from it}, and we name it only
because it ships with the code. A fifth, GDSC dose-response, is the external potency oracle used
for the Class-C analysis and is described below. Expression is processed to log-normalized
differential profiles against matched controls; per-context control means and treated-population
rows are exposed through a common dataset interface. Dataset scale, preprocessing settings and query
construction are tabulated in Supplementary Table 1.

\subsection*{DART retrieval task and distributional metrics}

Given a target population $Q$ and candidate populations $\{P_d\}$, DART ranks by a
population-to-population distance. Implemented scores: the energy distance\cite{EnergyDistance} (subsampled to a fixed maximum cell
count with a fixed seed), an RBF maximum-mean-discrepancy\cite{MMD}, a
sliced-Wasserstein distance (128 projections)\cite{SlicedWasserstein}, and a
subpopulation-coverage aggregate with a temperature parameter $\beta$
interpolating from mean aggregation ($\beta \to 0$) to worst-case
($\beta \to \infty$). Mean-signature retrieval, by cosine or by Euclidean distance between
mean signatures, is the zero-variance limit of the same family.

\subsection*{Signature and PCA baselines}

CMap-style signature retrieval is implemented with two selectable scoring rules. The
cosine rule computes the cosine of mean-delta signatures and is the same operation as
mean-cosine retrieval, which is the source of the exact Hit@1 identity. The WTCS rule is a
rank-based (GSEA-style) weighted connectivity score, a monotone sibling that is not
algebraically identical. PCA-latent retrieval fits PCA per query and scores in latent space
by mean or distributional distance.

\subsection*{Predict-then-rank baselines}

Candidate responses are generated by an average-effect predictor (per-drug mean
deltas), a nearest-neighbour predictor (nearest training context by control-state
similarity), and a CPA-style linear-latent predictor, used as a surrogate for
the latent-arithmetic family that scGen\cite{ref4} belongs to; we do not benchmark a trained
scGen, and the active backend is recorded in every output. They are used purely as
candidate-response \emph{sources}, not as models under benchmark. They do not stand in
for the broader class of published perturbation predictors and virtual-cell foundation models
(scGen\cite{ref4}, CPA\cite{ref5}, chemCPA\cite{ref6}, GEARS\cite{ref10}, scGPT\cite{ref7},
scFoundation\cite{ref8}, Geneformer\cite{ref9}), and no claim about that class is drawn from
them. An earlier draft asserted that they did stand in for it; that is inconsistent with the
algebraic scope of the Gate-1 result and with limitation 1, and it is withdrawn.

\textbf{How a predicted effect becomes a predicted population.} A distributional score
ranks \emph{populations}, so a predict-then-rank baseline must turn a predicted effect into
a population, and that synthesis step silently determines what structure the population can
possibly contain. We therefore state it explicitly. A predicted population is built by applying
the predicted effect to the query context's \emph{real control cells}, one cell at a time.
This is what the latent-arithmetic models in this class do (encode observed cells, shift
each cell's code, decode each cell), and it is the only construction under which a
predicted population can carry the control population's covariance and any multimodality it
has. The alternative and superficially natural construction, drawing each predicted
population as $\mathrm{control\ mean} + \delta + \mathrm{iid}\ \mathcal{N}(0,\sigma)$, must
not be used for structure diagnostics: that cloud is unimodal by construction (isotropy
index 0.998), so a $k$-means $k{=}2$ between-total variance ratio computed on it returns
$\approx$0.009 whatever the predictor learned. Any structure diagnostic evaluated on such a
synthesizer measures the synthesizer, not the predictor. Both constructions are implemented
behind an explicit flag so the difference can be reproduced.

\textbf{In-sample ceiling.} In the predict-then-rank experiment the predictors are fit on
the full SciPlex3 matrix and the query is composed of cells from two cell lines that remain
in the training set, so the nearest-neighbour predictor can in principle transfer a
signature it estimated from the very cells that compose the query. The reported
predict-then-rank numbers are therefore an in-sample ceiling, not a transfer
result. A leave-both-contexts-out option is provided, and we do not claim a leave-one-context-out guarantee for the reported run. The
direction of this bias favours the mean incumbent, so it does not manufacture the negative
result reported here; it bounds it.

\subsection*{Gate 1 measured within a context, and the calibration check that caught us}

\textbf{The theorem is a within-context statement.} An additive predictor adds one delta vector to
every cell \emph{of the query context}. Its two subpopulations therefore receive identical deltas
and $\cos(d_0,d_1)$ is exactly 1, but only if the two subpopulations \emph{are in that context}.
Our first Gate-1 experiment measured the cosine on an alpha-blended candidate population whose two
subpopulations are two different cell lines, and an additive predictor emits a different delta per
cell line, so what it measured was $\cos(\delta_{\mathrm{A549}}, \delta_{\mathrm{K562}})$. That is
the similarity of two contexts' deltas and is not a test of additivity.

\textbf{The calibration check is what exposed it, and it is asserted rather than inspected.} The
experiment scores predictors that are additive \emph{by construction}, whose value is fixed by the
algebra at 1.000. They returned 0.267 and 0.353. A control that is required to hit a known value and
misses it condemns the instrument, not the models, and the script refuses to emit numbers when it
fails. We recommend the practice generally: an experiment whose theory pins a control to an exact
value should assert that value, not merely report it.

\textbf{The corrected design} (analysis/predictors/gate1\_within\_context.py). For each context we
take up to 400 control cells and split them into two subpopulations by their own cell state (30-component
PCA, $k$-means with $k=2$; subpopulation sizes and silhouettes are reported, because a low induced
cosine would mean nothing if the within-context structure were not there). For each drug the
predictor is run \emph{separately on each subpopulation's control cells}, so the partition is true
by construction and no assignment step is involved:
\[
d_k = \mathrm{mean}\big(\mathrm{predicted\ cells\ of\ subpop\ }k\big) -
      \mathrm{mean}\big(\mathrm{control\ cells\ of\ subpop\ }k\big), \qquad
\cos(d_0, d_1).
\]
Population synthesis draws a permutation of the supplied cells, which preserves a mean exactly. The
real treated cells of the same context, partitioned by the same control-derived subpopulations, give
the reference value. Gate 2's assignment error is no longer entangled with this measurement; it is
measured separately (above).

\textbf{Two baselines, because a generative model does not return $C + \delta$.} A latent predictor
returns $P(C) + W\delta_z$, where $P$ is its own reconstruction of the input. The response of a
subpopulation therefore depends on what it is referred to, and the two choices are not the same
quantity:
\[
\underbrace{d_k = \mathrm{mean}\big(P(C_k)\big) - \mathrm{mean}(C_k) + W\delta_z}_{\text{against the
REAL control cells: what a scorer measures}}, \qquad
\underbrace{d_k = W\delta_z}_{\text{against the model's own decoded control}} .
\]
The first carries $\mathrm{mean}(P(C_k)) - \mathrm{mean}(C_k)$, the reconstruction bias of
subpopulation $k$, which is cell-state dependent and does not cancel; the two subpopulations' biases are
themselves near-orthogonal ($\cos = -0.05$ to $-0.18$). The second cancels it exactly.
We report both. The additivity theorem is a statement about the second, and the calibration check is
asserted on it: an additive predictor must return $1.000000$, and a run in which it does not is
discarded. A predictor that does not reconstruct its input (the gene-space average-effect predictor,
and the optimal-transport map, which returns $C$ plus a per-cell displacement) has no baseline of its
own, and for those two the columns coincide by construction, which is the correct answer rather than
a bug. Which predictors reconstruct and which do not is declared explicitly in the code rather than
inferred.

\subsection*{Information-condition diagnostics and gate diagnosis}

The information condition is tested as two gates. Gate 1 (differential response)
uses the predict-then-rank experiments: candidate responses are generated by each predictor
family, ranked by DART, and compared to mean retrieval, with structure diagnostics computed
on generated versus real populations. Two quantities are reported. The
\emph{subpopulation-variance ratio} is the $k$-means $k{=}2$ between-total variance ratio of
the candidate population. The \emph{induced response cosine} is
$\cos(d_{\mathrm{maj}}, d_{\mathrm{min}})$, where $d_k$ is the mean of subpopulation $k$
minus the mean of its own matched control population; it is the quantity a distributional
score must exploit, and it is identically zero for any predictor that adds a single delta
vector to every cell. The reference population for both is the alpha-blended candidate
population that a scorer actually ranks, not a single-context treated population.

\textbf{Gate 2} (recoverability) asks whether a query's subpopulation structure can be resolved at
all, on a controlled mixture with known two-state composition (K562 at 10\,$\mu$M, HDAC-class
versus JAK-class cells; 50--400 cells per source, 10 seeds). The design point is that clustering
failure on its own is uninformative: it is equally consistent with the information being absent
from the representation and with its being present but unreachable without labels, and those have
opposite implications for whether single-cell resolution can ever pay off. The analysis therefore
reports a supervised ceiling alongside the unsupervised methods, and the decisive quantity
is the \emph{gap} between them, not either number against chance.

\emph{Common unit.} Adjusted Rand index and classifier AUC are not comparable, so every method is
converted to best-permutation accuracy against the true labels: each cluster is assigned its
majority true class. For $k = 2$ this is the better of the two permutations; for $k > 2$ it is
strictly generous, since an over-partitioning method is never penalised for splitting a true class.
HDBSCAN noise is treated as its own cluster and likewise assigned its majority class, which is
generous again. Clusterers thus receive their label matching for free.

\emph{Unsupervised methods:} $k$-means and Gaussian mixtures, both handed the true $k = 2$; Leiden\cite{Leiden}
(graph-based, resolution 1.0, which selects its own $k$); HDBSCAN\cite{HDBSCAN} (density-based, which
selects its own $k$ and may return none). The last two are the field's actual defaults and were absent from our
earlier analysis; adding them does not change the answer. All four see the same PCA representation
(30 components) fit on the full mixture, which is the transductive setting an unsupervised pipeline
would in fact have.

\emph{Supervised ceiling:} logistic regression, random forest and $k$-nearest-neighbours, each
given the ground-truth labels under 5-fold stratified cross-validation, with the PCA
\emph{re-fit inside each training fold} and test cells only projected through it, so no test-fold
covariance informs the representation. The ceiling per run is the best of the three; including a
nonlinear learner is what distinguishes an information limit from a too-weak-probe artifact.

\emph{The direction of the resulting bias, corrected.} An earlier version of this Methods section
stated that giving the clusterers the easier transductive representation while making the probe
re-fit its PCA in-fold biased the comparison \emph{against} the conclusion drawn. That is
backwards. The conclusion drawn is that the \emph{gap} is small; inflating the unsupervised arm
and deflating the supervised arm makes the gap \emph{smaller}, so the design flatters
that conclusion rather than challenging it. We therefore also report the ceiling under the
\emph{matched} representation, the same transductive PCA the clusterers receive, which leaks no
labels because PCA does not see them, and it is the matched gap that should be read
(analysis/natural/gate2\_drug\_response.py).

\emph{Positive control.} Source separation is scaled by a factor $s \in \{1, 1.5, 2, 3, 5, 8\}$
using a transform fit once on the full source populations, never on the sampled cells; $s = 1$ is
the real data. Verifying that the ceiling pulls away from clustering as $s$ rises is what licenses
reading the absence of a gap at $s = 1$ as a measurement rather than as an insensitive probe.

Query identifiability on real drugs is the silhouette of a two-way $k$-means split computed per
drug on the query cells alone.

The gate is additionally audited directly by re-running the partial-observed retrieval with
an additively computed true response divergence column (leave-out-free, from the
observed candidate populations), then correlating each gate axis with true divergence and
stratifying retrieval gaps by divergence quartile with Benjamini--Hochberg control and power
analysis. This audit is the basis for the findings that the structure-reliability axis is
anti-correlated with true divergence and that the minority-coverage effect, while
significant, is negligible.

\textbf{Dataset positions on the divergence gate.} For each dataset the response cosine
between the two subpopulations of a query is computed from the observed populations with the
same probe used inside the gate, giving 0.032 (cross-line SciPlex3), 0.709 (Frangieh) and
0.186 (CD34$+$).

\subsection*{HIR-Bench generator and synthetic welfare oracle}

HIR-Bench generates heterogeneous retrieval tasks from a specified two-subpopulation latent
welfare model with tunable conflict, yielding an analytical flip boundary
$\alpha^* = B/(A+B)$. Each task is instantiated under the observed, predicted\_mean, and
predicted\_structure information conditions. Reported numbers come from the full
13,440-\emph{instance} grid (672 parameter configurations $\times$ 20 seeds; $n_{+} = 8{,}640$,
$n_{-} = 4{,}800$).

\textbf{The predictability model.} A logistic regression predicts retrieval flips from two
feature sets, fitted and reported side by side because the contrast between them is the
result. The \emph{observable} set contains only quantities a retrieval method can compute at
query time: the subpopulation-variance ratio, isotropy index and response diversity of the
query population and of the candidate populations. The \emph{oracle-derived} set
contains top-$k$ disagreement and Kendall conflict scores, each of which is a function of the
benchmark's latent utility matrix, as is the label.

Folds are held out one \emph{label-determining parameter cell} at a time, that is, one
$(\alpha, \mathrm{conflict})$ pair, because the label is a deterministic function of those
two axes; there are 28 such cells and this is the effective sample size. Grouping by the instance identifier instead, the natural choice, is invalid here: that
identifier also encodes axes that leave the utility matrix unchanged, so each held-out instance leaves
roughly twelve near-identical copies of itself in the training set. We report the AUC under
that grouping (0.640 for the oracle-derived features) only to quantify the inflation it
produces. Because the label is a deterministic step function of the grouping variables,
every held-out fold is single-class, no per-fold AUC is defined, and the pooled out-of-fold
AUC is the only available statistic.

\subsection*{Partial-observed retrieval protocol and real-data projection}

The partial-observed protocol scores candidates by an energy-based welfare regret against
the latent objective and records both the objective-aligned regret (Class A) and
task-proximal metrics (Class B: MoA-recovery nDCG\cite{nDCG}, minority-state coverage) per query.

\textbf{Held-out unit, and what it does not license.} The held-out unit in the gate audit
is a cell line (A549, K562, MCF7), and all three come from SciPlex3, so they share
batch, protocol, highly-variable-gene selection and normalization. This is
leave-one-cell-line-out within a single experiment, not leave-one-dataset-out, and we make
no cross-dataset generalization claim on its basis.

\textbf{Undefined metric values are excluded, not encoded.} MoA recovery is \emph{undefined}
for the leave-MoA-out and partial-library splits, where the entire mechanism class is removed
from the candidate library, so no correct answer is present to recover. Those 165 of 765
queries (21.6\%) are dropped, giving $n = 600$. They must not be carried as a sentinel value:
differencing a sentinel against itself yields a structural zero that enters every stratum as
a spurious no-effect observation and pins the median gap at exactly 0.000. Excluding them
does not rescue the metric; it makes the negative result sharper.

\textbf{Real-data projection.} The projection maps the four datasets onto the HIR-Bench
regime plane and records, per task, the minority-state coverage of the DART pick minus that
of the mean pick. The full configuration builds \textbf{239} tasks, and the composition is given here because the three-factor description in an earlier draft accounted for only 180 of them: 90 constructed cross-line tasks (three cross-line pairs $\times$ ten drugs $\times$ three seeds), 90 within-line (three cell lines $\times$ ten drugs $\times$ three seeds), 23 Frangieh, 12 CD34$+$ and 24 under predicted candidates. We report the distribution of that difference.

\textbf{Why the dominance count is withheld.} A task is conventionally called
distributionally dominant when the difference exceeds $+0.01$. We compute the count and do
not report it as a finding, for three reasons. (i) The threshold is not calibrated to the
data: it sits at 0.77 standard deviations of the nonzero-difference distribution
($\mathrm{s.d.} = 0.013$), and it is one-sided, so on the full run 11 tasks clear it for DART
while 2 clear it for the mean and only the first number would conventionally be quoted.
(ii) The seed is not a replicate. It re-draws which drugs are held out and it redefines the
minority subpopulation by re-clustering, so a task recorded as dominant at one seed carries
no information about robustness. (iii) When the drug is held fixed and the seed is resampled
ten times, only 2 of the 10 threshold-crossing cross-line tasks remain above threshold in at
least 8 of 10 seeds, while control drugs that cross it at no seed in the recorded run cross
it in \CONTROLRANGE{} of 10 resamples. A control here is a drug that is dominant at no
recorded seed; selecting controls by task identifier alone would let a drug that crosses at
one seed and misses at another fall into both groups. The count is a count of noise
excursions. The distribution is not, and it is what we report.

Two caveats bound the projection. First, the regime boundary it
transfers is fitted on HIR-Bench's method-performance layer, which is available on a
144-instance grid because scoring nine retrieval methods against up to 100 candidates per
instance does not scale to the full grid. The boundary is therefore fitted on 72 grid cells,
and we claim no more precision for it than that supports. Second, the projection's own
acceptance test fails. Its label derives from the same threshold we decline to trust, and it is heavily imbalanced
(228 to 11), so a constant predictor scores 95.4\% while the regime projection scores 61.1\%,
which is 34 points \emph{worse} than saying nothing. The HIR-Bench boundary does not transfer
to real data, and that holds whatever one makes of the threshold.

\subsection*{Evaluation metric taxonomy and circularity audit}

We classify every evaluation metric into three classes. Class A: objective-aligned
proxy metrics measure the same mathematical object the method optimizes or scores (energy-distance
regret for energy-based retrieval, connectivity score for signature retrieval, reconstruction
error for expression predictors, graph proximity for graph inverse design); they test
objective fidelity, not independent utility. Class B: task-proximal metrics are
biologically closer but still derived from the same upstream data (MoA recovery,
target-family recovery, DEG overlap, pathway enrichment, minority-state coverage).
Class C: oracle-independent utility metrics have external functional grounding
(dose-response viability, longitudinal survival, post-treatment clone expansion, validated
resistance emergence, functional readouts, combination efficacy) and are required for
therapeutic-utility claims.

A field-level audit across six published method families, spanning signature
retrieval\cite{ref1,ref2,ref3}, expression predictors and foundation
models\cite{ref4,ref5,ref6,ref7,ref8,ref9}, and graph/target inverse
design\cite{ref10,ref11,ref12}, finds that every family's \emph{headline} metric falls in Class A
and that none reports a Class-C readout as its primary evidence; population-level distributional
metrics are themselves standard in recent perturbation-prediction benchmarks\cite{ref15}. The
audit, its inclusion criteria, and the rule by which alignment risk is judged are given in full as
Supplementary Table 2, so the claim can be checked rather than taken on assertion. It is an audit
of \emph{evaluation practice}, not of correctness, and the family is the unit: we do not score
individual papers. We report DART under all three classes. No Class-C readout exists within the
single-cell datasets themselves; the Class-C analysis below therefore imports one, and we label it
semi-real for that reason.

\subsection*{Class-C analysis: an independent functional oracle and its scalar controls}

\textbf{Choosing the oracle.} A retrieval score is asked which candidate's response distribution
most resembles the query's. It is not asked which candidate is most cytotoxic. Grading it on
\emph{absolute potency} therefore poses a question the method was never set, and it does so in a
way that is not neutral: potency is largely driven by response magnitude, so a scalar that sorts
candidates by magnitude is near-guaranteed to win, and a similarity retriever handed a weak query
will correctly return other weak drugs and be scored as if it had failed. We therefore take the
\emph{semantically matched} oracle to be drug--drug functional similarity and report the
absolute-potency comparison only as a confounder audit.

\textbf{Construction.} SciPlex3 drugs were matched to GDSC records by normalized drug name (exact
match, substring matching as a fallback), giving 35 matched compounds after removing one
duplicate mapping (two SciPlex3 entries resolving to the same GDSC compound) and restricting to a
\emph{single} screening platform, GDSC2. That platform restriction matters: GDSC1 and GDSC2 AUCs
come from different assays and dose ranges and are not on a common scale, so pooling them would
inject an assay-of-origin step into the oracle itself. For each matched compound we take its
dose-response AUC profile across GDSC2 cell lines and define
\[
\mathrm{oracle}(q,c) \;=\; \mathrm{Spearman}\big(\mathrm{AUC}_{\bullet}(q),\;
\mathrm{AUC}_{\bullet}(c)\big),
\]
the rank correlation of the two compounds' cross-cell-line sensitivity patterns. Three properties
make this a usable Class-C oracle. It is \emph{query-dependent}, so a query-independent scalar
cannot game it. It is \emph{out of context}: the three SciPlex3 lines (A549, K-562, MCF7) are
\emph{excluded} from the profile, leaving \CTXLINES\ cell lines, so the oracle is never evaluated
in the biological context the transcriptional data came from. And because a correlation is
centered per compound, a drug's overall potency \emph{level} cannot drive it; only its selectivity
\emph{pattern} can, which is precisely the channel the absolute-potency comparison confounds.

Cell lines differ greatly in general drug sensitivity (mean AUC across all GDSC2 compounds ranges
from 0.720 to 0.966), and that shared main effect inflates every drug--drug correlation: left
uncorrected, 85\% of drug pairs appear positively similar (median $+0.178$). We therefore subtract
each cell line's mean AUC, estimated across all \NGDSCDRUGS\ GDSC2 compounds rather than only the
35 matched ones, before correlating; the corrected oracle is symmetric about zero (52\% positive,
median $+0.015$). The correction was specified on this reasoning before the retrieval results were
inspected. Both oracles are reported: they change every scorer's absolute value but not the
ordering that carries the conclusion. Pairs sharing fewer than 100 cell lines are dropped, leaving
591 usable drug pairs at a median of 909 shared lines.

\textbf{Task.} Per cell line, leave-one-drug-out: the held-out drug's 10\,$\mu$M treated
population is the query and the remaining matched drugs are the candidates (\NCQUERIES\ queries;
34--35 candidates per line). Each rule ranks the candidates and the ranking is compared with the
oracle by Spearman correlation. The oracle is visible to no scorer.

\textbf{The rankings and their controls.} (i) \emph{energy}, the distributional score. (ii)
\emph{mean cosine with control subtraction}, the correctly specified incumbent (treated mean minus
matched control mean). (iii) \emph{mean cosine without control subtraction}, how a mean baseline
is often run in practice, which measures baseline cell identity rather than drug response. (iv)
\emph{magnitude matching}: rank by $-|\,\|d_c\| - \|d_q\|\,|$, preferring candidates whose response
is about as large as the query's. This is query-dependent yet compares no distributions
whatsoever, and it is the null a distributional score must beat. (v) \emph{magnitude only}: rank by
$\|d_c\|$, which is query-independent and produces one ordering for every query; on a
query-dependent oracle it is a falsifiable sanity check rather than a competitor. (vi)
\emph{potency matching}: rank by $-|\mathrm{AUC}_c - \mathrm{AUC}_q|$ in the query's own line. This
is a \emph{diagnostic}, not a deployable recommender, since it requires the query drug's own
measured potency; it tests how far the functional oracle is re-encoding potency level, and we
report that it does so substantially (median $\rho = +0.399$).

\textbf{Reported statistics.} Per-query Spearman correlations summarized by their median. Energy is
additionally partialled on the magnitude-matching channel by rank-residualizing both the score and
the oracle on it, which isolates the contribution attributable to the distribution rather than to
magnitude. The \NCQUERIES\ queries within a cell line share a candidate pool, so paired tests
across them are anticonservative; we report such a $p$-value once, explicitly flagged, and treat
the three cell lines as the honest unit, giving per-line medians throughout. We do not convert any
rank correlation into a fraction of variance explained; a Spearman $\rho$ is not one.

\textbf{Confounder audit (absolute potency).} The same rankings scored against absolute GDSC2 AUC
are reported in Extended Data. There, energy is anti-correlated with potency ($\rho = -0.520$) and
the query-independent magnitude scalar reaches $+0.692$. This follows arithmetically: response
magnitude predicts potency ($\rho = -0.58$ to $-0.71$ by line), and the energy \emph{distance}
between query and candidate largely tracks the candidate's own magnitude ($\rho = +0.791$), so
ranking nearest-first returns the smallest-response candidates. We report this to expose the
magnitude channel, and we draw no conclusion about therapeutic utility from it.

\subsection*{The oracle-shape test: holding the data fixed and changing only the judge}

\textbf{Data.} Frangieh Perturb-CITE-seq measures RNA and 20 surface proteins (24 antibodies, 4
isotype controls removed) in the same cells. The two modalities are joined by cell
barcode, not by row position: the RNA and protein matrices as distributed differ in row count, so
a positional concatenation would pair each cell's protein response with a different cell's RNA,
which produces a result that looks unremarkable and is pure noise. The join is asserted, giving
218,331 cells. Protein counts are CLR-normalized across antibodies per cell. Analyses are run
\emph{per immune condition} (Control, IFN$\gamma$, Co-culture; these live in
\texttt{perturbation\_2}, not \texttt{condition}), because pooling them confounds condition with
perturbation. Knockouts with at least 60 cells are used, subsampled to 250, seed fixed.

\textbf{The two oracles.} Both are built from the \emph{same} protein measurements of the
\emph{same} cells, and differ only in statistical form:
\[
\mathrm{oracle}_{\mathrm{MEAN}}(q,c) = \cos\big(\bar{d}^{\,\mathrm{prot}}_q,\;
\bar{d}^{\,\mathrm{prot}}_c\big), \qquad
\mathrm{oracle}_{\mathrm{DIST}}(q,c) = -\,\mathcal{E}\big(D^{\mathrm{prot}}_q,\;
D^{\mathrm{prot}}_c\big),
\]
where $\bar{d}^{\,\mathrm{prot}}_k$ is the mean protein response of knockout $k$ against that
condition's control mean, $D^{\mathrm{prot}}_k$ is the set of \emph{per-cell} protein responses of
knockout $k$ against the same control mean, and $\mathcal{E}$ is the energy distance. Centering both
on the control is what makes the distributional oracle a \emph{response} distance rather than a
baseline-identity distance.

\textbf{What is held fixed.} The two RNA rankings under test, energy and control-subtracted mean
cosine, are computed once and scored against both oracles, so they are identical across the two
columns by construction. No scorer sees protein data of any kind. Nothing but the oracle's form
changes.

\textbf{The magnitude control, and why it is mandatory here.} Both distributional objects in this
design are energy distances, and an energy distance largely tracks a candidate's own response
magnitude ($\rho = +0.791$; Extended Data Fig.~5). A magnitude-to-magnitude channel could therefore
reproduce the reversal without any distribution being compared, and a paper that requires this
control of others cannot omit it from its own sharpest claim. We therefore add a third RNA ranking,
magnitude matching ($-\big|\,\|d_c\| - \|d_q\|\,\big|$), which is query-dependent and compares no
distributions, and score it against both oracles. Reported statistics are per-query Spearman
correlations summarized by their median, with the three immune conditions given separately as the
honest unit.

\subsection*{Resistance-associated exploratory analysis}

Response divergence per cell is correlated with curated resistance-associated program scores
(AXL/mesenchymal, interferon-response, antigen-presentation) in the Frangieh data. This
analysis is exploratory; no survival, timecourse, or functional readout is available, and no
resistance-prediction claim is made.

\subsection*{Natural heterogeneity: patient-derived glioblastoma}

\textbf{Data.} ZhaoSims2021: acute-slice culture and surgical biopsy from ten glioblastoma
patients, treated with six compounds (Ana-12, Ipsenisib, RO4929097, Tazemetostat, etoposide,
panobinostat). Counts are normalized to $10^4$, log1p-transformed, and reduced to 2{,}000 highly
variable genes. This is not a retrieval benchmark and is not used as one. Only one patient
received more than two compounds, so a leave-one-drug-out retrieval would run on a single patient
against five candidates; reporting Hit@1 from that would be an anecdote presented as a benchmark.
It is used to test the two gates, and the paper's motivating premise, on heterogeneity that nobody
constructed.

\textbf{Compartment assignment, and why it is not an argmax.} The published annotation carries
\texttt{celltype = unknown} for every cell, so compartments are assigned by canonical marker score:
each cell receives a $z$-scored mean expression for each marker set, computed on the full
log-normalized matrix before highly-variable-gene selection so the labelling never depends on a
gene set chosen for downstream variance. A first version of this procedure took the bare argmax
over five marker sets, which \emph{forces} every cell into some compartment including cells that
express no markers at all. It produced two compartments that fail their own validation: cells
labelled T cells (22\% of the dataset) had a mean CD3D/CD3E/CD2 expression of 0.080, \emph{lower}
than their malignant-marker score of 0.141, and cells labelled endothelial expressed myeloid
markers more strongly than endothelial ones. They were neither; they were unidentifiable cells
swept into whichever class the argmax had left over. Building a natural-heterogeneity result on a
fabricated T-cell compartment would be worse than having no natural dataset.

A cell is therefore assigned to a compartment only if it scores highest for that compartment,
clears an absolute expression floor on that compartment's key markers ($\geq 0.25$ mean
log-normalized expression), and beats the runner-up by a $z$-score margin ($\geq 0.25$). Everything
else is labelled unassigned and \emph{excluded}, not reassigned; 39.9\% of cells are dropped this
way. The natural-tissue result is insensitive to these two constants. Over
$\mathrm{floor}=\mathrm{margin}\in\{0.10,0.25,0.40,0.50\}$, which drops between 29\% and 58\% of
cells and passes compartment validation at every setting, the Gate-2 ceiling stays $0.915$--$0.923$,
the best unsupervised method $0.77$--$0.80$, the gap $0.11$--$0.13$, and the Gate-1 induced cosine
$0.54$--$0.59$ (analysis/natural/zhao\_threshold\_sensitivity.py). The 0.25 setting is neither the
most nor the least favourable. Only the three compartments that pass validation are used: malignant glioma cells (own-marker
expression 1.02 versus 0.07 and 0.12 for the others), tumour-associated myeloid cells (1.53 versus
0.12, 0.17) and oligodendrocytes (1.78 versus 0.41, 0.11). The loader raises rather than writing
labels if any called compartment fails to express its own key markers most strongly, and the
validation table ships with the results. Labels are never shown to any clustering method or to any
retrieval score; they exist only to evaluate whether those methods recover them.

\textbf{Gate 1 (natural).} For each (patient, drug) with at least 50 cells in both the malignant and
the myeloid compartment, $d_k = \mathrm{mean}(\text{treated cells of compartment }k) -
\mathrm{mean}(\text{that patient's control cells of compartment }k)$, and the reported quantity is
$\cos(d_{\mathrm{malignant}}, d_{\mathrm{myeloid}})$. Controls are matched \emph{within patient and
within compartment}, so the cosine cannot be inflated by a shared patient or compartment effect.

\textbf{Gate 2 (natural), and the comparison our first attempt got wrong.} Gate 2 asks whether the
subpopulations \emph{inside a candidate population} can be told apart, which is a
\emph{drug-response} partition. Our first natural test instead separated malignant from myeloid
\emph{control} cells, a \emph{cell-type} partition between two lineages that differ in thousands of
genes. It is a far easier problem, it returned a ceiling of 0.964, and comparing that with the
constructed mixture's 0.692 compared two different questions while holding the protocol fixed. That
number is withdrawn.

The reported analysis (analysis/natural/gate2\_drug\_response.py) poses the same question in both
settings. A split is a pair of drugs: cells treated with drug A against cells treated with drug B,
\emph{within one compartment and within one patient}, at least 50 cells per drug, capped at 150 per
source, 5 seeds. This gives 36 drug-pair splits over 4 patients and 2 compartments. The constructed
comparator is run at the same granularity, one drug against one drug in K562 at 10\,$\mu$M
(24 HDAC-versus-JAK drug pairs, subsampled with a fixed seed from the 513 available and reported as
subsampled). The class-pooled split the paper previously used is retained as a third arm, because a
pooled class carries the extra within-class variance of several distinct compounds and is therefore
harder for a reason that has nothing to do with the mixture being constructed. Quoting it against
the natural arm would inflate the contrast.

\emph{The representation, and a bias we had stated backwards.} Everything else follows the
constructed-mixture protocol: a common 30-component PCA for the unsupervised methods,
best-permutation accuracy as the shared unit with label matching given to each clusterer for free,
and a supervised ceiling taken as the best of logistic regression, random forest and
$k$-nearest-neighbours under 5-fold cross-validation. We report the ceiling under two
representations. The \emph{in-fold} ceiling re-fits PCA inside each training fold. This is the
leak-free choice, but it is not the representation the clusterers were given; because it lowers
the ceiling it \emph{shrinks} the gap, which is the direction that supports the conclusion the gap
is used to draw. The \emph{matched} ceiling uses the same transductive PCA the clusterers receive,
which leaks no label information because PCA never sees labels, and it is the apples-to-apples
quantity; it is what we report. \emph{The reported gap is the median of the paired per-split
differences}, not the difference of the two marginal medians; a median is not linear and the two
disagree (for the natural arm, $+0.117$ against $+0.147$), and the paired statistic is the correct
and the smaller one.

\textbf{Uncertainty on the gap} (analysis/natural/gate2\_uncertainty.py). Three things make the raw
number fragile and each is addressed. (i) The gap is a difference of two \emph{maxima}, best of three
probes minus best of four clusterers, and a maximum over methods is upward-biased. We therefore also
compute every fixed (probe, clusterer) pair, in which no selection occurs; HDBSCAN is excluded from
the summary because it assigns every cell to noise and scores exactly $0.500$, so a gap measured
against it reflects a clusterer that did nothing rather than information that is unreachable.
(ii) The confidence interval is a cluster bootstrap over patients, not over splits. Thirty
of the 36 natural splits come from one patient and reuse the same six drugs, so a split-level
bootstrap would treat dependent observations as independent and return an interval several times too
narrow, which is precisely the pseudo-replication this paper criticises elsewhere. With only four
clusters the bootstrap is coarse and is known to under-cover, so the interval is reported as a floor
on the uncertainty rather than a faithful estimate of it. (iii) The dominant patient is dropped
entirely and the analysis repeated on the remaining three. The resulting values are given in Fig.~\ref{fig:5}d, with two that appear only here: the maximisation moves the natural number by only $-0.005$, and the natural-minus-constructed difference is $+0.109$ $[+0.052, +0.129]$, with none of 5{,}000 bootstrap draws at or below zero.

\textbf{The premise.} For every pair of drugs within a patient, we compute the cosine between their
population \emph{mean} signatures and the cosine between their \emph{malignant-compartment}
responses, and report the Spearman correlation between the two across drug pairs. If the mean were a
sufficient statistic these would be the same number. Because the mean signature contains the
malignant cells, 43\% of those called, we also report the \emph{disjoint} form, in which the
malignant-compartment response is ranked by the myeloid-compartment response instead; the two
compartments share no cells and each is referred to its own compartment-matched control
(analysis/natural/zhao\_premise\_disjoint.py). \emph{No pass/fail threshold is applied.} An
earlier version of the analysis script declared the premise supported if the correlation fell below
0.9, and the observed value was 0.878. That threshold was ours, it was applied after seeing the
number, and clearing it by 0.02 licenses nothing; the verdict has been removed from the code and
the correlation is reported and interpreted in the text.

\subsection*{Identifiability phase diagram}

The phase diagram sweeps a separation scale and a cell budget on the K562 HDAC-versus-JAK
mixture, whose labels are known by construction. Separation is applied by pushing each source
along the vector from its centroid to the global centroid. Those centroids are computed once,
on the \emph{full} source populations, so the manipulation is independent of the sample being
evaluated. This is not a detail: computing the centroids on the sampled cells and then
shifting those same cells fits the manipulation on the data it is measured on, amplifies
whatever separation each small draw happens to contain, and produces the indefensible result
that the adjusted Rand index \emph{falls} as the cell budget rises. Clustering does not become
harder with more cells.

\subsection*{Statistical analysis}

Paired comparisons use the Wilcoxon signed-rank test; group comparisons use Mann--Whitney U;
monotone associations use Spearman correlation. Multiple comparisons across strata use
Benjamini--Hochberg false-discovery-rate control\cite{BenjaminiHochberg}, and we report achieved statistical power and
the sample size required for 80\% power where effects are small. Effect sizes are reported
alongside $p$-values throughout, and statistically significant but negligible effects are
labeled as such. Where an effect is reported as negligible, the operating range of the metric
is given so that the effect size can be read against it.

%% =====================================================================
%%  BACK MATTER
%% =====================================================================
\section*{Data availability}
\sloppy

All datasets analysed in this study are publicly available.

Three of the four single-cell datasets were obtained in harmonized form from
scPerturb\cite{ref3}, Zenodo record 13350497 (DOI
\href{https://doi.org/10.5281/zenodo.13350497}{10.5281/zenodo.13350497}, release 1.4), and we name
the exact file for each so the analysis can be reproduced from the same bytes rather than from a
re-derivation of the primary data. SciPlex3: file
\texttt{SrivatsanTrapnell2020\_sciplex3.h5ad}; primary data in Gene Expression Omnibus under
accession GSE139944\cite{dataSciPlex3}. Frangieh melanoma Perturb-CITE-seq: files
\texttt{FrangiehIzar2021\_RNA.h5ad} and \texttt{FrangiehIzar2021\_protein.h5ad} (the protein
modality is used by the surface-protein oracle); original data\cite{dataFrangieh}.
ZhaoSims2021 patient-derived glioblastoma (the natural-heterogeneity benchmark): file
\texttt{ZhaoSims2021.h5ad}; original data\cite{dataZhao}. CD34$+$ HSPC: Gene Expression
Omnibus, accession GSE306429.

\textbf{GDSC dose-response} (the Class-C functional oracle): Genomics of Drug Sensitivity in
Cancer, GDSC2 fitted dose-response, release 8.5 (file
\texttt{GDSC2\_fitted\_dose\_response\_27Oct23.xlsx}), available from
\url{https://www.cancerrxgene.org} and the Sanger bulk-download mirror. The specific release was
recovered and verified: the retained join table (\texttt{drug\_match\_table.csv}) reproduces the
AUC and $\ln \mathrm{IC}_{50}$ of this release exactly across all checked (compound, cell line)
pairs, and the Class-C statistics were re-derived from it to confirm they are unchanged. The join
table is released with the code so the drug matching can be audited independently.

\textbf{LINCS L1000 / PDGrapher benchmark}: a precomputed benchmark table ships with the code
(\texttt{pdgrapher\_closed\_loop\_benchmark.parquet}). No reported result in this paper derives
from it, so it is not data underlying the findings; we name it only because it is present in the
release, and we do not have a retained accession for it.

The processed tensors used by the experiments, together with the scripts that build them from the
primary data, are described in the repository. Source data for every figure panel are provided as
the experiment output tables in the code repository.

\section*{Code availability}

All code, experiment scripts, and the evaluation-metric taxonomy are available at
\url{https://github.com/Boom5426/DART} (\TODO{archived DOI on Zenodo at acceptance}). The
release includes a test suite and a numerical consistency check that re-derives every
reported constant from the cached experiment outputs; the latter re-reads the cached
experiment tables and diffs them against the recorded constants, and is a consistency check
rather than an end-to-end re-execution. A single driver script re-executes the pipeline from the processed tensors.

\section*{Acknowledgements}
\TODO{Acknowledgements. Funding sources and grant numbers.}

\section*{Author contributions}
\TODO{Author contributions, CRediT taxonomy: conceptualization, methodology, software,
formal analysis, investigation, data curation, writing (original draft), writing (review and
editing), visualization, supervision, funding acquisition.}

\section*{Competing interests}
\TODO{Competing interests statement. If none: ``The authors declare no competing interests.''}

\section*{Additional information}
\noindent Supplementary Information is available for this paper: Extended Data Figs.\
1--6, Supplementary Notes 1--3 and Supplementary Tables 1--4.\\
Correspondence and requests for materials should be addressed to \TODO{Name}.

\bibliographystyle{unsrtnat}
\bibliography{references}

\end{document}
